\section{五次の高分母：局所分類と大域整合}
\label{nt19:denominators}

この節には，今後の高分母の構成・排除に再利用する条件を残す．
高分母を持つ全平面五次の大域作用 atlas の存在・不存在は，なお一般には未解決である．
現在の必要候補は五頂点全てを使う五型・六十六標識であり，
最小次数二以上，二点連結で，高辺は七本以上である．
二価頂点の二つの高近傍はともに次数四であり，その二価頂点を含む高分母三角形を作る．
二価頂点数は零・一・三のいずれかで，七高辺の場合は book だけが残る．
高分母が存在するなら最大値は百三以上で，少なくとも二種類の高分母が現れる．
この有界排除と，分母上限を用いない graph 障害とは量化が異なる．

三高分母の正符号 graph は完全 matching であることが証明されたため，
固定低データの二次復元と canonical Schur 線形系を全ての正整数三高模型に使える．
三高係数の同時非負性は，$b=-\beta\ge2$，$d=-D_0\ge1$，
$\alpha\ge k\ge b+1$ の非特異領域に限られる．
第21版では五つの整数 $(s,t,D_a,D_b,D_c)$ を固定した全候補の有限復元を得た．
この五変数そのものの一様有限化と，残る五型の五 core 同時接着は未解決である．
Book の残余二領域と $G_{5,15}$ の有限中心還元を，以下の具体的な入口として使う．

\subsection{正値性と二高分母の表示}

局所四頂点に $q_i>0$，$p_{ij}\in\ZZ$ を置く．
\[
 \sum_{j\ne i}p_{ij}=-q_i,\qquad
 p_{ij}p_{ji}=q_iq_j-1,\qquad
 D_i=-q_i-\sum_{j\ne i}p_{ji}.
\]
正の basin weights の存在は，$p_{ij}>0$ の directed edges からなる全 simple cycles で
\[
 \prod_{i\to j}p_{ij}<\prod_iq_i
\]
となることと同値である．二閉路は determinant 条件で自動的に通るので，
最大八個の有向三角形と六個の有向四角形を検査すればよい．
これは既知の差分不等式・max-times の判定\cite[\S2]{SergeevSchneiderButkovic}である．
全 gains に有理数 $\theta>1$ を十分一に近く掛け，終点 $j$ の simple paths の gain と一の最大を
$w_j$ と取れば，正の有理 weights の証明書を構成できる．

二高分母 $(a,b,1,1)$，$a,b>1$ では，high vertices および low vertices の交換を許すと，
次の表示が全 numerator matrices を尽くす．
\[
 r>0,\quad r\mid a-1,\quad r+1\mid b-1,\quad
 u\ne0,\quad u\mid a-1,\quad
 x=\frac{a-1}{r},\quad y=\frac{b-1}{r+1},
\]
\[
 A=a+x+u,\quad A\mid ab-1,\quad
 B=\frac{ab-1}{A},\quad v=B-b+y\ne0,\quad v\mid b-1.
\]
このとき $A,B>0$ であり，
\[
 P=\begin{pmatrix}
 0&-A&x&u\\
 -B&0&-y&v\\
 r&-r-1&0&0\\
 (a-1)/u&(b-1)/v&d&0
 \end{pmatrix},\qquad
 d=-1-\frac{a-1}{u}-\frac{b-1}{v}.
\]
正の weights が存在するための必要十分条件は
\[
 u<0<v
\]
と分かっている．算術稿の \texttt{ls11:two-high-signs} は $a=b$ を仮定しない．
従って正の模型を探索するとき，$u>0$ の候補を残す必要はない．
各 high row のもう一つの high への entry は負であり，二つの low への entries は一正一負となる．
固定した $(a,b)$ の局所幾何は，この有限な divisor 計算と局所 chart の十分条件で判定できる．
分母対 $(a,b)$ の全体が有限個という主張ではない．

奇数 $b>1$ に対し，$(3,b,1,1)$ の正の局所 chart が存在することと
$b\in\{7,11\}$ は同値である．この結論にも $b$ の上限はない．
非正の numerators だけなら全 $q_i=1$ なので，通常の arithmetical structures の
非正非対角行列\cite[\S\S1--3]{CorralesValencia}へ高分母全体を帰着しない．

\paragraph{一般二高分母の負係数と，再利用する整数変数．}
正の標準形で $u=-s<0$，$v=z>0$，$S=(a-1)/s$，$T=(b-1)/z$ と置くと，
$d=S-T-1<0$ であり，$\min(D_A,D_B)<0$ も一般に成立する．
証明は算術稿の \texttt{td12:low-entry} と \texttt{td12:one-negative} にある．
二高分母が同時に非負の $D$ を持つ候補を再探索する必要はない．
再利用する変数は
\[
 h=S-r,\quad k=T-r-1,\quad
 \alpha=x-s=sh/r,\quad \beta=y-z=zk/(r+1)
\]
であり，これらは全て非零で同符号となる．さらに
\[
 M=rs+1,\quad N=(r+1)z+1,\quad
 \delta=\alpha-\beta,\quad L=M-N+r\delta
\]
と置けば
\[
 M\beta-N\alpha=1,\qquad \alpha L=1+a\delta,
 \qquad D_A=h-L,\quad D_B=L+\delta-k
\]
となる．特に $\gcd(\alpha,\beta)=1$ である．
負枝 $\alpha,\beta<0$，$\delta\geq0$ では，整数性が $S=T$ を強制する．
この枝では $d=-1$ かつ $D_B<0$ となる．

正の局所 weights は
\[
 (w_A,w_B,w_C,w_D)
 =\left(a,b,x+\frac1{2r},z+\frac1{2T}\right)
\]
と明示できる．正の entries に対応する四つの basin areas は
$1/(2r),1/2,1/(2T),1/2$ となるため，整数データから面積正値性の証明書を直接作れる．
この定理は分母の上限，奇数性，大域 clutching，node budget を用いない．
ただし一般の三高分母で全ての高係数 $D$ が非負となることを禁止できるかは未解決である．
三高分母の numerators の matching 性とは別の問題である．

\paragraph{再試行を避けるための符号反例．}
一意の最大分母の係数が負とは限らない．二高分母だけでも
\[
 q=(19,17,1,1),\qquad
 P=\begin{pmatrix}
 0&-23&6&-2\\-14&0&-4&1\\3&-4&0&0\\-9&16&-8&0
 \end{pmatrix},\qquad
 w=\left(19,17,\frac{37}{6},\frac{33}{32}\right)
\]
は正の局所模型であり，$D=(1,-6,5,0)$ となる．
従って閉路に沿う $D$ の向きから分母の単調性を無条件には結論しない．
三高分母の既存例 $(13,7,7,1)$ でも $D=(4,-2,-4,2)$ となる．
また，異なる二高分母の係数和も常に負ではない．
具体的には
\[
 q=(11,1,1,9),\qquad
 P=\begin{pmatrix}
 0&-2&5&-14\\
 -5&0&0&4\\
 2&-1&0&-2\\
 -7&2&-4&0
 \end{pmatrix},\qquad
 w=(33/16,1,1,33/8)
\]
は正の局所 chart であり，$D=(-1,0,-2,3)$ を与える．従って
\[
 \sum_{q_i>1}D_i=2>0,\qquad
 \sum_i\frac{D_i}{q_i}=-\frac{58}{33}<0.
\]
一方，等二高分母 $q=7$ の標準模型では $D=(0,-2,2,0)$ であり，
$\sum_iD_i/q_i=12/7>0$ となる．したがって，単純な高分母部分和や
この重み付き和の一様な符号を，大域の potential として仮定しない．

\paragraph{等三高分母の負和．}
分母 $(q,q,q,1)$，$q\geq3$，では低分母頂点の係数は $D_0>0$ であり，
高分母三頂点の係数和は負となる．この結果には正値性を用いない．
正局所模型では一般の matching 定理からも，低分母 row が二負一正であることが分かる．
その row を $(-a,-b,a+b-1)$ と書けば
\[
 D_0=-1+(q-1)\left(\frac1a+\frac1b-\frac1{a+b-1}\right)>0
\]
となる．正値性を仮定しない元の等三高分母定理は，別途 row/pair の整数条件から証明されている．
等二高分母の解決済みの分類と，それに使った整数変数を次にまとめる．

% v11 reusable notes. The equal-two-high inequality and integer classification are solved.
\subsection{等二高分母の完全分類と再利用する整数変数}
\label{uh10:two-reduction-section}

等二高分母 $(q,q,1,1)$ の局所符号問題は解決した．
高分母二頂点の係数和は常に負であり，正の整数 numerator data 自体も完全に分類できる．
主論文の \texttt{uh11:two-negative} と \texttt{uc11:classification} に一般証明を置く．
ここでは，他の分母 pattern に移す際に使える変数と，分類の結果だけを残す．

\paragraph{整数性を保つ表示．}
一般二高分母の符号定理から，標準表示では $u=-s<0<v=z$ である．
\[
 m=q-1=t r(r+1),\qquad R=r(r+1),\qquad r,t\geq1,
 \qquad s,z\mid m,\quad s,z>0
\]
と置くと，残る条件は
\begin{equation}\label{nt11:two-conic}
 sz+(1+tr^2)s-\{1+t(r+1)^2\}z=t+1
\end{equation}
となる．低分母の二頂点の係数は
\[
 D_C=\frac mz-\frac ms-t,\qquad D_D=s-z-1.
\]
従って，この全整数解について解決した不等式は
\[
 D_C+D_D=(s-z)\left(1+\frac m{sz}\right)-t-1>0,
 \qquad D_A+D_B=-(D_C+D_D)<0
\]
である．等しい分母という仮定を外した係数和には，同じ符号を仮定しない．

\paragraph{一般の上限を作る補助整数．}
\[
 S=m/s,\quad Z=m/z,\quad H=Z-S,\quad
 \delta=s-z,\quad n=2r+1,
\]
\[
 1\leq\delta\leq t,\qquad e=t-\delta,\qquad
 G=r^2H+R-nS=\frac{R(e+1)}{sz}>0,\qquad J=H-\delta G,
\]
\[
 tJ=(\delta-1)H,\qquad Jsz=\delta(\delta-1)R,
 \qquad\delta\mid Jz,
\]
\[
 K=nz-r^2\delta+1>0,
 \qquad tK=sz+\delta-1
\]
を同時に保つ．$G,J,K$ を単なる実数へ緩めると，証明に使う情報を失う．
特に $\delta\geq2$ では $x=z/\delta\geq r^2/n$ と
$Jx(x+1)=(\delta-1)R/\delta$ から $J\leq4$ を得る．
$J=4$ は $4x\in\ZZ$ によって排除され，$J\in\{1,2,3\}$ となる．
さらに一般の係数評価で $J\leq\delta-2$ を排除するため，
\[
 \delta-1\leq J\leq3,\qquad \delta\leq4
\]
となる．その証明の最後の四つの二次方程式は，一般の不等式から残った候補であり，
分母を有界にした探索の結果ではない．

$\delta=1$ では $s=z+1$，$R=HK$，$sz=tK$ となる．
$a=r-z>0$ と置き，$a^2=(t-1)K+1$ と $(2r+1)a=(H-1)K+1$ を消去すると
\[
 LK=t+2H+1,\qquad L=(H+1)^2-4Ht>0
\]
を得る．従って $H=4t-2+k$，$k\geq0$ と書け，$K$ の整数性が $k=0$，
または $k=t=1$ だけを残す．この消去は，有限の例外と無限族を分ける際に再利用できる．

\paragraph{全整数データと分母集合．}
標準表示の $(r,t,s,z)$ は，標識の交換を除いて次表で尽きる．
最終行では $t$ は任意の正整数である．
\begin{center}
\small
\begin{tabular}{ccccccl}
$r$ & $t$ & $s$ & $z$ & $q$ & $\delta$ & $D_C+D_D$ \\
\hline
$2$ & $1$ & $2$ & $1$ & $7$ & $1$ & $2$ \\
$1$ & $6$ & $4$ & $1$ & $13$ & $3$ & $5$ \\
$2$ & $5$ & $6$ & $3$ & $31$ & $3$ & $2$ \\
$6t-3$ & $t$ & $3t$ & $3t-1$ & $t(6t-3)(6t-2)+1$ & $1$ & $3t-2$
\end{tabular}
\end{center}
従って，可能な等高分母はちょうど
\[
 \{7,13,31\}\ \cup\
 \{t(6t-3)(6t-2)+1:t\geq1\}
\]
である．無限族の最初の分母は $13,181,721,1849,\ldots$ となる．
$q=13$ の第二行と無限族の $t=1$ は別の整数模型であり，
high--high pair の絶対値はそれぞれ $(21,8)$ と $(14,12)$ である．
この分類は numerator systems の分類であり，weights の変形や幾何学的同型類の数ではない．

無限族では，高分母二頂点の係数は標準表示の順で
\[
 D_A=2(3t-2)(t-1),\qquad D_B=-(3t-2)(2t-1).
\]
$t\geq2$ なら一方は正であるが，総和は常に $-(3t-2)<0$ である．
大域排除に使えるのはこの総和であり，個々の高分母端の係数が常に負という主張ではない．

\paragraph{現在の限界と次の利用先．}
等二高分母の再探索や，上の不等式の反例探しは必要ない．
次は異なる分母 $(a,b,1,1)$ の端点接続で，この無限族がどの denominator や
opposite $D$ と整合し得るかを調べる．単色成分は既に大域に存在できないため，
局所等高分母模型を用いる候補には，その同じ高分母成分内に異なる値も必要となる．


\subsection{単色の高分母成分と二値分母の排除}

高分母 support を $H=\{ij:q_{ij}>1\}\subset K_5$ とする．
辺を持つ一つの連結成分 $C$ の全高分母が同じ $q>1$ なら，
各 core の分母 pattern は $(q^k,1^{4-k})$ となる．
$k=1,2,3$ では高分母側の $D$ の和は負であり，$k=4$ では全四頂点の和なので零である．
しかし $C$ 上で足すと各 edge の対向端が相殺し，総和は零でなければならない．
従って全 vertices の次数が四となり，$C=K_5$ となるが，これは素数 support の
$K_4$ 禁止に反する．

従って，高分母 support の各非空連結成分には少なくとも二種類の高分母が必要である．
特に，ある $q>1$ に対して全 branch denominators が $\{1,q\}$ に属すれば，全て一となる．
算術稿の \texttt{uh11:monochromatic-component} と
\texttt{uh11:two-value-rigidity} に対応する．
この障害は分母上限，node 非負性，全 node budget を用いない．
「単色の部分 graph」が別の高分母 edge に接する場合には，連結成分という仮定が失われるので，
その部分 graph だけを同じ議論で排除しない．

\subsection{素数 support と共通因子の制約}

五 cores の labels を $V=\{0,1,2,3,4\}$ とする．
共有 branch の分母を $q_{ij}=q_{ji}$，core $i$ の outgoing numerator を $p^i_{jk}$ と書く．
\[
 \sum_{k\ne i,j}p^i_{jk}=-q_{ij},\qquad
 p^i_{jk}p^i_{kj}=q_{ij}q_{ik}-1,\qquad
 q_{ij}\mid p^i_{jk}+p^j_{ik}.
\]
奇素数 $\ell$ ごとに $G_\ell=\{ij:\ell\mid q_{ij}\}\subset K_5$ とする．
その三角形 $ijk$ 上では，$\mathbf F_\ell$ において
\[
 p^i_{jk}=p^j_{ki}=p^k_{ij}=\varepsilon,\qquad
 p^i_{kj}=p^j_{ik}=p^k_{ji}=-\varepsilon,
 \qquad\varepsilon=\pm1.
\]
三つの determinants と三つの clutchings を一周すると $x^2=1$ を得る．
この式と row sums から $K_4$ と book $B_3$ が禁止され，$G_3$ は matching となる．

現在の密な support に対する一般結論は
\[
 |E(G_\ell)|\geq7
 \quad\Longrightarrow\quad
 \ell=5,\quad G_\ell\cong K_{2,2,1}
\]
である．従って exact 七辺 support はどの奇素数でも存在せず，
\[
 |E(G_3)|\leq2,\qquad |E(G_5)|\leq8,\qquad
 |E(G_\ell)|\leq6\quad(\ell\ne3,5).
\]
五の場合も edge 数七は不可能である．七辺 graph の残存二型
$K_1\vee P_4$ と $K_{2,3}$ に一本を加えた graph は，
全 cores の式を同時に使うと八本目の辺を強制する．
後者では最後の determinant の残差が $-5$ になることが，素数を五へ限定する核心である．

四 cores に対する以前の gcd 条件に加え，任意に選んだ七本の異なる branches に対して
\[
 \gcd\{q_{ij}:i,j\in S,\ i<j\}=1\quad(|S|=4),
 \qquad
 \gcd(q_{e_1},\ldots,q_{e_7})\mid5
\]
となる．特に二十五で割れる分母は高々六本である．
$G_5=K_{2,2,1}$ の場合，その八辺の分母の gcd は正確に五であり，
少なくとも二辺では五進付値が一となる．
この指数の障害は $\ZZ/25\ZZ$ 上で証明する．
その環では $r^2=0$ から $r=0$ としてはならず，nilpotent ideal を保持して
最後の determinant 残差 $-5\ne0$ を取り出す．

固定された標識付き $K_{2,2,1}$ に対し，modulo five の reduced system は符号 $s=\pm1$ の二解だけである．
parts を $\{a,b\},\{c,d\},\{o\}$ とすると
\[
 q_{ab}=-s,\quad q_{cd}=s,\quad q_{ij}=0\ (ij\ne ab,cd),
 \qquad
 P^o=s\begin{pmatrix}
 0&-2&1&1\\-2&0&1&1\\-1&-1&0&2\\-1&-1&2&0
 \end{pmatrix}
\]
となり，他の四 cores も一意に決まる．これらは対向 $D$ の相殺も満たすため，
同じ mod~$5$ 条件だけでこの型をさらに排除できない．
正の奇整数への lifting，全整数 clutching，面積，node 配分，actual neck の接着は別に確認する必要がある．

三に関する局所計算も再利用できる．一つの core で二つの分母が三で割れると仮定すると，
それらの rows は $\mathbf F_3$ 内で一定符号となり，残る二 rows の最初の二 entries は打ち消す．
最後の pair は $p_{23}=-q_2$，$p_{32}=-q_3$ を要求し，determinant の $-1$ と矛盾する．
実際の atlas で全ての高分母が三で割れるなら，$H=G_3$ が matching となり，
孤立高分母 edge の障害から $H$ は空である．
素数 support の合同定理には weights や node 非負性を使わないが，
この最後の幾何的結論には正の basin areas と対向 $D$ の相殺を使う．

\subsection{低分母 puncture と三角形成分の一般障害}

同じ高分母 $q$ の共有 branch の一つの marked puncture が，両側で低分母一へ向かうとする．
その outgoing entries $p,p'$ は $q-1$ の非零 signed divisors であり，
$q\mid p+p'$ から
\[
 p'=-p,\qquad\{p,p'\}=\{1,q-1\},\qquad
 \{p,p'\}=\{-1,-(q-1)\}
\]
の三型に限られる．各型で incoming entries の和も $p+p'$ に等しい．
meridian node 表示での個数は順に $1,0,2$ である．
これは marked punctures の対応を固定して使う必要がある．

$H$ の一つの連結成分が三角形なら，各 core は二高分母を持つ．
二つの外側 low punctures の outgoing sums は，符号が揃う場合に $q,-q$，
逆に対応する場合に零，零となるため，総和は常に零である．
従って三角形に向かう high puncture の outgoing numerators の和は $-2q$ となり，
対向 $D$ の相殺も使うと high--high numerators の差は一つの circulation $\delta$ に一致する．
それらの絶対値の平均を $T_i>0$ と書けば
\[
 q_i=\frac{T_i+T_{i+1}}2,\qquad
 (T_{i-1}+T_i)(T_i+T_{i+1})=4T_i^2-\delta^2+4.
\]
最大値と最小値の比較から $\delta^2=4$ と全 $T_i$ の一致を得る．
従って三辺分母が等しくなり，単色成分の障害に反する．
局所の追加計算では $q=7,13$ の high--high pair $(-q-1,-q+1)$ だけが残り，
その高分母係数は $(0,-2)$ または $(0,-1)$ で，やはり相殺できない．

この三角形成分の排除は分母の上限によらず，node 非負性や全 node 数五も仮定しない．
三角形に別の高分母 edge が接続する場合は，この議論の範囲に含まれない．
一般の path 成分には，次に述べる別の局所符号の議論を用いる．

\subsection{固定した低 row と column からの三高分母の完全復元}
\label{nc17:fixed-low}

今後は三高分母 core の全分類を待たず，隣接 core から低データが決まった時点で
次の二次復元を使える．低 ray を $0$ とし，整数
\[
 r_i=p_{0i},\quad x_i=p_{i0},\quad
 \sum_i r_i=-1,\quad q_i=1+r_ix_i>1\quad(i=1,2,3)
\]
を固定する．三つの $r_i$ の同符号の対を $r_2,r_3$ とし，
$s_i=-q_i-x_i<0$，$Q_{ij}=q_iq_j-1>0$ と置く．
$A=p_{12}$ は
\[
 (s_2A-Q_{12})\{s_3(s_1-A)-Q_{13}\}
 -Q_{23}A(s_1-A)=0
\]
を満たす．二次係数 $a=Q_{23}-s_2s_3$ は，$r_2,r_3>0$ なら負，
$r_2,r_3<0$ なら正なので消えない．従って高高 entries の completion は高々二つである．
その整数根から
\[
 p_{21}=Q_{12}/A,\quad p_{13}=s_1-A,\quad p_{31}=Q_{13}/(s_1-A),
\]
\[
 p_{23}=s_2-Q_{12}/A,\qquad p_{32}=s_3-Q_{13}/(s_1-A)
\]
を復元する．必要十分な整数条件は $A\ne0,s_1$，$A\mid Q_{12}$，
$s_1-A\mid Q_{13}$ であり，その後に全正 cycle の条件を課す．
判別式は
\[
 \Delta=(s_1s_2s_3-s_1Q_{23}-s_2Q_{13}-s_3Q_{12})^2
            -4Q_{12}Q_{13}Q_{23}.
\]
任意の探索上限を置かず，整数平方根と二つの割り算だけで全候補を得る．

実装 \nolinkurl{research/exact_three_high_completion_v17.py} の
\texttt{complete\_three\_high} は元の物理標識の順を保存する．
同符号対を選ぶ内部の標識交換を，隣接 core への clutching に誤って持ち込まない．
通常の整数行列だけを求める場合と，正の weights を求める場合を区別できる．
\nolinkurl{verify_three_high_completion_v17.py} は別の signed-divisor 復元との一致も検査する．

分母 $q$ と低 row $r$ が既知なら，$x_i=(q_i-1)/r_i$ も決まるため同じ手順を使える．
一方，分母を可変のまま低 row だけを固定した場合には高々二つとはいえない．
既存の rank 二 Schur 無限族は低 column 自体が変化する例である．
$r=(-2,-2,3)$，$x=(-3,-3,2)$，$q=(1,7,7,7)$ は二つの正の標識付きcompletionを持つが，
両者は高 ray $1,2$ の交換で対応する．標識付き解の個数を幾何学的同値類数と同一視しない．

% Reusable finite bound; optional research-note material, not needed for K4 exclusion.
\subsection{反対符号の二つの例外 reverse が大きい場合の有限化}
\label{fc17:tail-note}

低 row と column が
\[
 r=(P,-P+h-1,-h),\qquad x=(1,-1,-v),\qquad
 P\ge h\ge1,\quad v\ge1
\]
である整数三高分母模型を考える．分母は
\[
 q=(P+1,P-h+2,hv+1)
\]
となる．高高 entry $t=p_{13}$ を用いて他の entries を行和・pair式から消去すると，
\[
 P\{vt^2+2(hv+1)t+(h-1)(hv+1)\}
 +vt^2+\{v(h+1)+1-h\}t+hv(h-1)=0
\]
を得る．$t\ne0$ であり，$t>0$ なら左辺は正なので，$t=-z$，$z\ge1$ でなければならない．
\[
 \begin{split}
 A(z)&=vz^2-2(hv+1)z+(h-1)(hv+1),\\
 B(z)&=vz^2-\{v(h+1)+1-h\}z+hv(h-1)
 \end{split}
\]
と置けば $PA(z)+B(z)=0$ であり，
\[
 B(z)-A(z)=\{(h-1)v+h+1\}z-(h-1)>0.
\]
従って $A(z)<0$ で，
\[
 P+1=\frac{\{(h-1)v+h+1\}z-(h-1)}{-A(z)}.
\]
$A(z)=B(z)=0$ という共鳴は生じない．

$z\ge2h+2$ なら
\[
 A(2h+2)=(v-1)(h^2+3h+4)+(h^2-1)\ge0
\]
であり，その先で $A$ は増加する．また
\[
 B(z)=vz(z-h-1)+(h-1)z+hv(h-1)>0
\]
なので矛盾する．従って
\[
 1\le z\le2h+1,\qquad
 P\le\{(h-1)v+h+1\}(2h+1)-h.
\]
よって $h,v$ を固定すれば全 $P$ を完全に復元できる．
全分母が奇数なら $P$ は偶数，$h$ は奇数である．
通常 reverse の上界 $h\le4$ が使える場合には $h=1,3$ のみとなる．
この有限化は他の graph で反対符号の大きい例外 entries が残る場合に転用できる．


\subsection{高分母 support の現在地と，適用仮定}
\label{ns19:current-support}

対象は，全平面五次の標準 $K_5$ branch graph と $K_4$ cones を保つ作用適合 atlas である．
$H=\{ij:q_{ij}>1\}$ とする．木・単純閉路成分，二価二頂点を持つ三角形，
葉辺付き四角形，$K_{2,3}$，diamond，二葉付き三角形，四使用頂点上の $K_4$ に加え，
$G_{5,5},G_{5,10},G_{5,11},G_{5,12}$ と葉付き $G_{5,7},G_{5,9}$ を排除した．これらの障害は分母の数値上限を用いない．
現在の必要候補は\emph{五型・六十六標識}であり，整数接着や atlas の存在を主張しない．

非空 $H$ は五頂点全てを使う二点連結 graph で，最小次数二以上，少なくとも七辺を持つ．
全ての二価頂点は高分母三角形の中にあり，その二つの高近傍はともに次数四である．
二価頂点数は零・一・三に限られる．一個の場合は $G_{5,15}$，三個の場合は book $G_{5,13}$ である．
高辺が七本なら book であり，他の四型は八辺以上を持つ．
一つの core に高分母 ray が高々一本なら全分母一である．
これは現在の graph に対する結論であり，別の core や branch singularity を導入する場合には，
変更後の接着条件を定義してから障害の適用範囲を調べる．

\begin{center}
\begin{tabular}{lcc}
障害 & 対向 $D$ の相殺 & node 非負性\\\hline
木成分・三角形成分 & 用いる & 不要\\
四角形・五角形成分 & 不要 & 用いる\\
二価二頂点を持つ三角形 & 不要 & 用いる\\
四角形と一本の葉辺 & 用いる & 用いる\\
$K_{2,3}$・負の枝の三角形条件 & 不要 & 用いる\\
diamond・二葉付き三角形の全排除 & 用いる & 用いる\\
全 support が $K_4$ & 低 hub 辺で用いる & 用いる\\
$G_{5,5}$ & 二本の低辺で用いる & 用いる\\
$G_{5,10}$ & 用いる & 用いる\\
$G_{5,12}=K_4+\mathrm{leaf}$ & 葉の高辺で用いる & 用いる\\
$G_{5,7}$ & 葉の高辺で用いる & 用いる\\
$G_{5,9}$ & 葉に接する四辺で用いる & 用いる\\
$G_{5,11}$ & 低辺で用いる & 用いる
\end{tabular}
\end{center}
node 全数五および各 node の上限五は row sums と非負性から導く．
三角形の任意の追加高辺を許したわけではなく，二頂点が二価という仮定を維持する．
以下には，今後も使える定量的な伝播式を残す．

\subsection{三高分母の完全マッチング定理：今後は仮定でなく入力として使う}
\label{nm18:matching}

整数行和・pair products を満たし，$q_0=1,q_1,q_2,q_3>1$ である局所模型を考える．
正の entries の無向 graph は，完全マッチングまたは低 ray を含む三角形と孤立頂点の和に限られる．
さらに，正の basin weights を持つことと，完全マッチング型であることは同値である．
分母上限，clutching，node，対向 $D$ 条件はいずれも不要である．
証明は算術稿の \texttt{mh18:matching} に置く．

再利用時には，低 row の三つの entries がちょうど二負一正であることを最初に課す．
正の低 entry の相手を $c$，負の二つの相手を $a,b$ とすると，
正符号 graph は $0c\sqcup ab$ であり，既存の canonical Schur 表示が全ての正の三高模型を覆う．
$\beta\le-1$ と，境界 $\beta=-1$ の完全な約数表示をそのまま使える．
ただし，全ての高 $D$ が非負となる可能性はまだ解決していない．
この係数問題の残りは $b=-\beta\ge2$，$d=-D_0\ge1$，
$\alpha\ge k\ge b+1$ の非特異領域である．
特異 Schur 系は全高 $D$ 非負の候補から除かれたが，
一般の正整数局所模型の分類にはその実在する無限族を保持する．

\paragraph{他の問題にも使える二つの証明操作．}
二つの道型では，同じ残余二次式
\[
 (kC-b)L^2+\{kd-aC+B(kC-b)\}L+d(kB-a)=0
\]
に対し，全係数の符号がそろう．正の根 $L$ がないことが，無上界の排除を与える．
三角形型では，正の三閉路の不等式を
\[
 (tS-H)K<1,\qquad (sT-K)H<1
\]
と書くと，整数性から $H=tS+\alpha,K=sT+\beta$，$\alpha,\beta\ge0$ となる．
残る pair equations を整理した
\[
 bS=AT+\alpha\beta,\qquad
 (A-1)U=CS+k(\alpha+1),
\]
\[
 (b+1)U=(C+1)T+(k+1)(\beta+1)
\]
は $A>1,b>0,U>T$ を強制する．$\delta=U-T\ge1$ に対して
\[
 A\{s(1-\delta)-k\beta\}=bsU+bk(s+t)+Ct\beta
\]
の左辺は非正，右辺は正となる．
ここでは $A=s\alpha-t$，$b=s-t\beta$，$C=s+(s+t+1)k$ であり，
この排除に用いる変数は canonical matching 型の変数と混同しない．
実数での $U>T$ を整数の $\delta\ge1$ へ強める一歩が不可欠である．

\paragraph{正値性を省く誤りを避ける検査例．}
分母 $(1,5,19,31)$ と
\[
 P=\begin{pmatrix}
 0&1&1&-3\\4&0&2&-11\\18&47&0&-84\\-10&-14&-7&0
 \end{pmatrix}
\]
は整数行和・pair products を全て満たし，正符号 graph は三角形である．
しかし閉路 $0\to2\to1\to0$ の gain は $188/95>1$ であり，正 weights はない．
行和・pair products だけから完全マッチングを結論してはいけない．

\nolinkurl{verify_three_high_matching_v18.py} は十一の多項式恒等式を係数ごとに確認し，
道型の係数の符号を正変数の平行移動で検査する．
併せて固定低データの小さい入力を既存の完全復元器へ通し，実装の整合性を点検する．
この入力格子の検査を，一般定理の証明として引用しない．


\paragraph{別の接着問題にも再利用する葉の伝播則．}
葉 $i$，二価 core $j$，その先の core $k$ を考える．両方の高辺を実際に接続すると，
ある $h\geq1$ により
\[
 a=q_{ij}=h(h+1)+1,\qquad b=q_{jk}=(h+1)(h+2)+1,
\]
\[
 D^j_i=a,\qquad D^j_k=-b,\qquad
 p^k_{ji}=-(h+1),\qquad p^k_{ij}=-(h+2)
\]
に強制される．特に $b>a$，$\gcd(a,b)=1$ である．
最初の分母対は $(3,7),(7,13),(13,21),(21,31)$ となる．
この結果は $k$ の高分母次数を仮定せず，node 条件も用いないため，
cycle を含む成分に付いた長さ二の pendant path にも適用できる．
現行の五候補では葉自体が存在しないが，この式は node 条件を緩めた接着問題にも使える．
ただしこの必要条件だけで接着可能性を結論しない．

第三 core の条件を省くと，二つの追加族
\[
 (a,b,D^j_i,D^j_k)=(2t+1,6t+1,3t,1-8t),
\]
\[
 (a,b,D^j_i,D^j_k)=(6t+1,20t-1,7t,3-18t)
\]
が残る．後者は $D^j_k>-b$ なので，一般の二高分母の局所不等式に移し替えてはならない．
ところが次の core から元の葉へ向かう entry は $b-1$ の signed divisor で，
前者では $3t+2\pmod{6t+1}$，後者では $10t+2\pmod{20t-1}$ を満たす必要がある．
前者は $t=1$ だけ，後者は全て不可能となる．
一本の辺での整合と，次の辺までの整合の差を示す再利用可能な例である．

\paragraph{木の中心で使う係数和の上界．}
三高分母 $(a,b,c,1)$ の分母一の vertex を $0$ とし，
$r_i=p_{0i},x_i=p_{i0}$ と置く．整数性と局所式だけで
\[
 \sum r_i=-1,\qquad r_ix_i=q_i-1>0,\qquad
 \sum_{i=a,b,c}D_i=1+\sum x_i
\]
である．ある $x_j\leq-1$ が存在し，他の $x_i\leq q_i-1$ なので
\[
 D_a+D_b+D_c\leq a+b+c-\min\{a,b,c\}-2.
\]
一高分母の葉は反対側に $D\geq q$ を要求し，上記の二価 core を経由しても同じ下界が届く．
従って三価の star とその一辺を細分した木は，中心でこの上界に反する．
四価 star は全 $D$ の和零で排除される．
正局所模型では完全 matching 定理により，low row とその逆向き列は二負一正である．
正の entry が高 ray $c$ に対応すれば $x_a,x_b\le-1$，$x_c\le q_c-1$ なので，さらに
\[
 D_a+D_b+D_c=1+x_a+x_b+x_c\le q_c-2\le\max\{a,b,c\}-2
\]
を得る．これは新しい matching 性から直ちに得る上界であり，
一般三高分母の「全高 $D>0$ 不可」を仮定するものではない．

\paragraph{四角形と五角形で再利用する node の使い方．}
二高分母の負の高低 entry は $-(q-1)$ にはならず，絶対値は $(q-1)/2$ 以下である．
従って，共通の低 puncture で両端が負となることは整数 clutching と矛盾する．
四角形では正の pairing が全周で同じ向きになる．mixed node の二型を伝播させると，
一方の型は第五の全低分母 core により canonical $S=T=2$ を強制する．
この局所型は $(r,s,z,a,b)=(4,4,5,9,11)$ だけなので，四つの core の共有分母が一致しない．
他方の型は逆向き約数 $(4,3,2,1)$ と真約数性が矛盾する．
この改訂証明は $D$ を使わない．

五角形では，正の pairing の vertex は高々一つである．残る四つの連続した負の pairing から，
三本の高辺の共通低 puncture は全て正となり，それぞれの両端の値は $\{1,q-1\}$ となる．
canonical な $A$ だけが正の最大約数 $q-1$ を持ち得るため，これらの端が伝播する．
その結果，ある低辺の node 数が $n=2-q<0$ となる．
この三本の例外辺の伝播は，整数 clutching と node 条件を同時に課す手順として再利用できる．

\subsection{二高分母の最小 entry と，誤った強化への反例}

canonical な非零 low--low entry $d=S-T-1$ は，$\alpha,\beta<0$ のときちょうど $-1$，
$\alpha,\beta>0$ なら $d\leq-3$ となり，$d=-2$ はない．
$d=-1$ では大きい高分母の $D$ が負，小さい方が正である．
ただし，大域標識で固定した向きの low--low entry は零の場合もある．

この最小 entry 型にも，正の高係数和を持つ無限族がある．$H\geq2$ として
\[
 r=2H,\quad S=T=H,\quad s=2H,\quad z=(2H+1)(H-1)
\]
と置くと，
\[
 a=2H^2+1,\qquad b=H(H-1)(2H+1)+1,
\]
\[
 D_A=2H^3-2H^2-2H-1>0,\qquad
 D_B=-2H^3+3H^2+H+1<0
\]
であり，$D_A+D_B=H(H-1)>0$ となる．
全 row sums と pair products は恒等式であり，明示 weights も存在する．
従って，高係数和の一様な負性を再び仮定しない．

\subsection{固定分母の相手と，小分母の係数を保持した伝播}

奇分母 $n>1$ を含む二高分母模型の相手集合 $\mathcal P(n)$ は，$n\ne7,13$ なら有限である．
七と十三だけが共鳴し，それぞれ $6k+1$ と $12k+1$，$k\geq1$，という無限族を含む．
これら以外の相手は有限個である．canonical な $A/B$ の役割を分けることが不可欠である．
$A$ 側を固定すれば
\[
 z\{A(r+1)-(1+rs)T\}=(r+1)(s-x-1),
\]
$B$ 側なら
\[
 s\{(1+(r+1)z)S-rB\}=r(y-z-1)
\]
を用いる．右辺が零の枝を落とさずに扱う．
\nolinkurl{verify_fixed_denominator_partners_v12.py} の
\texttt{partners\_for\_fixed(n)} は，全非共鳴標準形と共鳴枝を返す．

相手の分母だけでなく固定 ray の $D$ を保存すると，二高分母 core 同士の接続について
次の必要条件を得る．共有分母三・五は不可能である．七では，非零係数なら $\pm7$ で，
隣接高分母は $\{3,13\}$ または $\{3,19\}$，零なら両側とも七である．
十三では，非零係数なら隣接高分母 $\{5,25\}$ と係数 $\pm15$，または
$\{7,21\}$ と係数 $\pm13$ に限られ，零なら両側の隣接分母は各々十三または十七である．
この零係数の場合を一般の接続から落としてはいけない．
固定七・十三の共鳴族の係数 $7k-9,13k-14$ は，もう一つの二高分母 core により相殺できない．
全四十七非共鳴模型と二共鳴のデータは v13 の小分母伝播検算に保存した．

% v13 research notes: reusable local classification, not an unresolved C5 claim.
\subsection{二つの負の reverse divisors が四以下の完全な局所 catalog}
\label{nr13:bounded-reverse}

二高分母標準形の負の高低 entries は $-s,-y$ であり，
それぞれの reverse numerators の絶対値は $S,r+1$ である．
これらがともに四以下という局所条件を課すと，
分母の上限を置かずに全ての整数模型を有限表と三つの族に分けられる．
この catalog は，今後別の support に現れる共通 $\lambda=0$ punctures を調べる際に再利用できる．
五角形成分自体は v13 で排除されており，その残存候補を意味しない．

\paragraph{仮定と完全性．}
正の整数二高分母標準形で $a,b$ は奇数とし，
\[
 2\leq S\leq4,\qquad 1\leq r\leq3
\]
を仮定する．$S\geq2$ は負の低分母 entry の真約数性から従う．
このとき，次の表の二十二模型と三つの無限族だけが可能である．
表は canonical data の分類であって，標識を忘れた幾何的同型類を数えるものではない．

\begin{center}
\small
\begin{tabular}{rrrrrrrrr}
$r$&$S$&$T$&$s$&$z$&$a$&$b$&$D_A$&$D_B$\\\hline
1&2&5&1&2&3&11&6&$-10$\\
1&2&5&2&6&5&31&18&$-27$\\
1&3&8&2&5&7&41&21&$-36$\\
1&3&9&2&2&7&19&7&$-15$\\
1&3&9&4&6&13&55&23&$-41$\\
1&3&9&6&26&19&235&127&$-211$\\
1&3&11&8&2&25&23&$-9$&9\\
1&3&12&4&1&13&13&$-3$&$-2$\\
1&4&9&1&4&5&37&21&$-36$\\
1&4&10&1&2&5&21&11&$-21$\\
1&4&12&1&1&5&13&6&$-15$\\
1&4&12&3&5&13&61&26&$-49$\\
1&4&12&4&13&17&157&78&$-138$\\
1&4&13&9&28&37&365&177&$-312$\\
1&4&13&10&62&41&807&428&$-747$\\
1&4&15&6&2&25&31&$-4$&$-1$\\
1&4&15&25&4&101&61&$-63$&102\\
1&4&16&3&1&13&17&0&$-9$\\
1&4&16&13&2&53&33&$-31$&45\\
1&4&18&7&1&29&19&$-15$&15\\
2&3&5&2&6&7&31&21&$-25$\\
2&4&7&3&12&13&85&58&$-73$
\end{tabular}
\end{center}

無限族は $t,k\geq1$ を整数として
\[
\begin{array}{c|c|c|c}
 (r,S,T,s,z)&(a,b)&(D_A,D_B)&\text{parameter}\\\hline
 (1,2,6,2t-1,t)&(4t-1,6t+1)&(3,-7)&t\geq1\\
 (2,3,6,2k,1)&(6k+1,7)&(6-6k,7k-9)&k\geq1\\
 (3,4,6,3k,2)&(12k+1,13)&(12-12k,13k-14)&k\geq1
\end{array}
\]
となる．第一族では二つの分母がともに増大する．
後の二族は固定分母七・十三の共鳴であり，公開稿の
\texttt{sp13:no-resonant-continuation} により，固定 ray をもう一つの二高分母 core に
接続して対向 $D$ を相殺することはできない．

\paragraph{導出：小さい reverse data から有限個の factorization へ．}
$h=S-r$，$k=T-r-1$ と置く．$h=0$ は canonical 補助変数 $\alpha$ の非零性に反する．
$h<0$ の可能性は $(r,S)=(3,2)$ だけである．
負の枝の二分により $T=S=2$ となり，整除性から $s=3u,z=2v$，$u,v\geq1$ と書ける．
このとき $\alpha=-u,\beta=-v$ であり，Bézout 式は
\[
 (9u+1)(-v)-(8v+1)(-u)=1,
 \quad\text{すなわち}\quad u(1-v)-v=1
\]
となるが，左辺は負である．従って六組
\[
 (r,S)=(1,2),(1,3),(1,4),(2,3),(2,4),(3,4)
\]
だけを調べればよい．全て $h,k>0$ である．

\[
 R=r(r+1),\quad U=(r+1)(S-r),\quad V=r(T-r-1),
 \quad E=r^2T-(r+1)^2S+R
\]
と定義すると，高高 pair product は
\begin{equation}\label{nr13:bilinear}
 Esz+Vz-Us=R
\end{equation}
と同値である．$s,z\geq1$ より
\begin{equation}\label{nr13:T-bound}
 s(E-U)\leq R-V
\end{equation}
である．一方，正の枝の不等式 $rT>(r+1)S$ を保持する．
これらから $T$ は次の二十六候補に限られる：
\[
\begin{array}{c|c}
 (r,S)&T\\\hline
 (1,2)&5,6,7\\
 (1,3)&7,8,9,10,11,12,13\\
 (1,4)&9,10,11,12,13,14,15,16,17,18,19\\
 (2,3)&5,6\\
 (2,4)&7,8\\
 (3,4)&6
\end{array}
\]
実際，$r=1$ で $T\geq6S-4$ なら $E-U\geq0$，$R-V=4-T<0$ となり矛盾する．
$(r,S)=(2,3)$ で $T\geq7$，$(2,4)$ で $T\geq9$ も同様である．
$(3,4)$ で $T\geq7$ なら
$s(E-U)\geq9T-56>24-3T=R-V$ となって矛盾する．

$E\ne0$ のとき，式~\eqref{nr13:bilinear} を
\begin{equation}\label{nr13:factorization}
 (Es+V)(Ez-U)=ER-UV
\end{equation}
と書く．$C=ER-UV\ne0$ なら $d=Es+V>0$ は $|C|$ の正約数であり，
\[
 s=\frac{d-V}{E},\qquad z=\frac{Us+R}{d}
\]
を調べればよい．正整数性，$r\mid sS$，$r+1\mid zT$，$a=sS+1,b=zT+1$ の奇数性を課すと，
表示した二十二模型がちょうど残る．候補生成と棄却理由は
\texttt{bounded\_reverse\_catalog\_verification\_v13.json} に全て記録している．

$E=0$ の三候補は $r=1,T=4S-2$ であり，式は
$(S-1)(2z-s)=1$ となる．従って $S=2,s=2z-1$ だけが残り，第一族を得る．
$E\ne0,C=0$ の二候補は $(r,S,T)=(2,3,6),(3,4,6)$ であり，
式~\eqref{nr13:factorization} と $Es+V>0$ から，それぞれ $z=1,z=2$ となる．
残る整除性が $s=2k,s=3k$ を与え，後の二族を得る．
全ての模型は元の pair products と row sums を満たし，
標準形の明示 weights により正である．

\paragraph{この分類の射程．}
有限表の最大分母八〇七は完全な factorization の結果として現れた数であり，
計算時の探索上限ではない．ただし reverse divisors の上限四は仮定である．
他の問題へ適用するときは，その上限が本当に node 条件等から従うことを先に示す必要がある．
また局所 charts の表は，整数 marked clutching，node 配分，actual areas の大域的実現を保証しない．

% v20: complete matching parameterization, with simultaneous nonnegative high D reduced to a nonsingular sector.
% Singular positive charts remain relevant to the broader local classification.
\subsection{三高分母の Schur 補行列と，符号原理における整数性}
\label{nsch18:three-high}

ここでは $q_0=1$，$q_1,q_2,q_3>1$ とし，
$A=\operatorname{diag}(q_i)+(p_{ij})$ が
\[
 A\mathbf1=0,\qquad q_iq_j-p_{ij}p_{ji}=1
\]
を満たす局所問題を扱う．目標は，整数模型において
$D_i=-q_i-\sum_{j\ne i}p_{ji}$，$i=1,2,3$，の少なくとも一つが負になるか，
少なくとも三つ全てが狭義正となることはないかを決めることである．
本項ではいずれも一般定理とは主張しない．

\paragraph{仮定を追加しない一般帰着．}
$r_i=p_{0i}$，$x_i=p_{i0}$ とすると
\[
 \sum_{i=1}^3r_i=-1,\qquad q_i=1+r_ix_i.
\]
全ての $r_i,x_i$ は非零であり，各対 $r_i,x_i$ は同符号である．
分母一の diagonal entry を pivot とした Schur 補行列
\[
 B=A_{\{1,2,3\},\{1,2,3\}}-xr^{\mathsf T}
\]
は整数行列で，
\begin{equation}\label{nsch18:unit-diagonal}
 B_{ii}=1,\qquad B\mathbf1=0.
\end{equation}
実際，diagonal は $q_i-r_ix_i=1$ であり，row sums は
$-x_i-x_i\sum r_j=0$ である．従って，ある整数 $\alpha,\beta,\gamma$ により
\[
 B=\begin{pmatrix}
 1&\alpha&-1-\alpha\\
 \beta&1&-1-\beta\\
 \gamma&-1-\gamma&1
 \end{pmatrix}
\]
と表せる．残る三つの pair equations は
\begin{equation}\label{nsch18:linear-x}
 x_i(r_i-B_{ji}r_j)+x_j(r_j-B_{ij}r_i)=B_{ij}B_{ji}
 \qquad(1\leq i<j\leq3)
\end{equation}
であり，$B,r$ を固定すれば $x$ に関する線形方程式である．
これは $p_{ij}=B_{ij}+x_ir_j$ を元の pair equation へ代入して得られる．
また，column sums を取ると
\begin{equation}\label{nsch18:defect}
 D_0=-1-\sum_i x_i,\qquad
 D_i=-\sum_{j=1}^3B_{ji}+r_iD_0.
\end{equation}
従って $D_0>0$ なら $\sum_{i=1}^3D_i=-D_0<0$ であり，
高分母側の少なくとも一つは直ちに負となる．
さらに，正の weights を持つ場合は全 $D$ ベクトルが零である整数模型は存在せず，
$D_0=0$ でも高側に負の係数が現れる．
従って正の weights の下で残る場合は $D_0<0$ に限られる（\ref{nhd20:reduction}）．

この行列問題は，二次の principal minors が全て一である非対称行列の特殊な場合である．
Lin--Sturmfels\cite[\S2, Proposition~4]{LinSturmfels} は principal minors と cycle sums の
整数多項式による変換を与える．本問題ではさらに row sums，整数性，正の weights を課すため，
一般の principal-minor relations だけから係数の符号や整除性は結論しない．

\paragraph{全ての正整数三高分母模型を尽くす canonical 表示．}
完全 matching 定理（\ref{nm18:matching}）により，正の numerators の辺は
高 rays 間の $ab$ と，残る高 ray $c$ と低 ray $0$ を結ぶ $c0$ に限られる．
従って以下の表示は追加の matching 仮定を必要とせず，全ての正整数三高分母模型を扱う．
これは有限個の模型への分類ではなく，残る整数条件を共通の線形系に表す帰着である．

$s,t,S,T,U\ge1$，$u=s+t-1$ と置く．適切に $a,b$ を交換すれば $k\ge0$ を選べ，
\[
 (q_a,q_b,q_c,q_0)=(1+sS,1+tT,1+uU,1),
\]
\[
 P=\begin{pmatrix}
 0&tS+\alpha&-(uS+\alpha+1)&-S\\
 sT+\beta&0&-(uT+\beta+1)&-T\\
 -(sU-k)&-(tU+k+1)&0&U\\
 -s&-t&u&0
 \end{pmatrix}
\]
となる．$H=tS+\alpha,K=sT+\beta,L=uS+\alpha+1,N=uT+\beta+1,
M=sU-k,O=tU+k+1$ の六つが全て正であることを保持する．
$\alpha,\beta$ は整数であり，
$\mathcal A=s\alpha-t,\mathcal B=t\beta-s,\mathcal C=s+uk$ と置く．
残る三つの pair equations は
\[
 \begin{pmatrix}
  \mathcal B&\mathcal A&0\\
  -\mathcal C&0&\mathcal A+1\\
  0&\mathcal C-1&\mathcal B+1
 \end{pmatrix}
 \begin{pmatrix}S\\T\\U\end{pmatrix}
 =\begin{pmatrix}
 -\alpha\beta\\k(\alpha+1)\\-(\beta+1)(k+1)
 \end{pmatrix}
\]
にちょうど同値である．正の weights は
$w=(q_a,H+1/(2K),u+1/(2U),1)$ と明示できる．
幾何的な奇数分母条件は $sS,tT,uU$ が全て偶数であることとして保持する．
公開算術稿の \texttt{sh14:matching-setup} と \texttt{sh14:linear-reduction} は，
この表示と線形化の証明箇所である．matching 定理により全正整数三高模型へ適用できる．

この canonical 表示では
\[
 \mathcal A=s\alpha-t\geq1,\qquad
 \mathcal B=t\beta-s\leq-2,\qquad
 \mathcal C=s+uk\geq2,\qquad \beta\leq-1.
\]
$d=U+1-S-T=-D_0$ と置くと
\[
 D_a=sd-\beta-k-1,\quad D_b=td-\alpha+k,\quad
 D_c=\alpha+\beta+1-ud.
\]
従って全ての高係数が非負となるには
\[
 sd\geq\beta+k+1,\quad td\geq\alpha-k,\quad
 ud\leq\alpha+\beta+1
\]
が必要十分である．全て狭義正なら右辺をそれぞれ一だけ強める．
境界 $\beta=-1$ では $k=0,D_0>0$ が証明され，その範囲から全高係数非負の模型は生じない．
全ての正整数三高分母模型のうち，高係数符号の問題で今後調べる範囲は
$b=-\beta\ge2$，$d\ge1$，$\alpha\ge k\ge b+1$，$\Delta>0$ である．
全ての特異系はこの符号問題から除かれる（\ref{nhd20:reduction}）．
matching 性が解決したことと，この $D$ 符号の問題が解決したことを混同しない．

固定 $s\geq2,t\geq1$ に対する $\beta=-1$ の全解は，
\[
 K_0=s^2+st+t^2-s-t,\quad g=\gcd(s-1,t),\quad G\mid K_0g,\quad G>0
\]
を列挙し，$\alpha=(G+K_0)/(st)$ の正整数性と
\[
 S=\frac{\alpha(s-1)(s\alpha-t+1)}G,\quad
 T=\frac{\alpha s(s+t-1)}G,\quad
 U=\frac{\alpha s(s-1)}G
\]
の正整数性を課して得る．幾何的な奇数条件は $sS,tT,uU$ が全て偶数であることとして最後に課す．
この完全な有限手続きは，固定 $s,t$ に関するものであり，一般分母の有限性は主張しない．

\paragraph{特異系は実在する：分類から除外しない．}
係数行列の determinant は
\[
 \Delta=-\mathcal B(\mathcal A+1)(\mathcal C-1)
          +\mathcal A\mathcal C(\mathcal B+1).
\]
$\Delta\ne0$ では Cramer 分子の整除性を確認する．$\Delta=0$ では拡大行列のrankと整数解格子を調べる．
実際，$(s,t,\alpha,\beta,k)=(3,12,12,-4,1)$ では rank が二で，正整数全解は
\[
 (S,T,U)=(136+200n,291+425n,93+136n),\qquad n\geq0
\]
で尽くされる．分母は $(409+600n,3493+5100n,1303+1904n,1)$，
高係数は $(-997-1467n,-4007-5868n,4671+6846n)$ である．
この固定sliceは解決済みであり，同じ補助データの探索を繰り返さない．
全ての特異系の分類と，大域の五 core 接着は未解決である．
ただし特異系の高係数が全て非負になる可能性は排除済みであり，その候補として再探索しない．
正整数三高分母の非matching型は存在しないことが証明されており，再探索の対象ではない．

\paragraph{高係数の同時非負性は非特異領域だけに残る．}
\label{nhd20:reduction}
canonical 表示で $b=-\beta$，$d=-D_0=U+1-S-T$ と置く．
整数 row sums，pair products，正の weights のみを仮定すると，
\[
 D_a,D_b,D_c\ge0
 \quad\Longrightarrow\quad
 b\ge2,\quad d\ge1,\quad \alpha\ge k\ge b+1,
\]
\[
 \Delta=\mathcal C(\mathcal A-\mathcal B)
          +\mathcal B(\mathcal A+1)
 \ge(-\mathcal B)(-\mathcal B-1)>0.
\]
分母の上限，奇数性，clutching，node 条件は不要である．
したがって，階数二を含む全ての特異 Schur 系から
全高 $D$ 非負の模型が生じる可能性はない．
この結論は，特異な正整数局所模型そのものを排除するものではない．

再利用する証明の要点を残す．
\[
 D_a=sd+b-k-1,\quad
 D_b=td-\alpha+k,\quad
 D_c=\alpha-b+1-ud,\qquad
 D_a+D_b+D_c=d
\]
である．まず全 $D$ 零を除く．その場合は
$d=0,\alpha=k=b-1$ で，$b\ge2$ である．
\[
 m=b^2-b+1,\qquad K=s^2+st+t^2-s-t,\qquad
 r=(-s,-t,u)^{\mathsf T},\quad x=(-S,-T,U)^{\mathsf T},
\]
\[
 B_0=\begin{pmatrix}
 1&b-1&-b\\
 -b&1&b-1\\
 b-1&-b&1
 \end{pmatrix}
\]
と置くと，Cramer 恒等式は $\Delta=mK$，$K=m-1$，
$x=r-B_0r/m$ を与える．整数性から各成分で $m\mid B_0r$ である一方，
\[
 B_0^{\mathsf T}B_0=m(3I-J),\qquad
 \lVert r\rVert^2=2m-1,\qquad \mathbf1^{\mathsf T}r=-1
\]
である．ここで $J$ は全成分が一の三次行列である．従って
$m^2\mid\lVert B_0r\rVert^2=m(6m-4)$，従って $m\mid4$ となる．
$m$ は三以上の奇整数だから矛盾する．
特に $D_0=0$ なら高係数には正と負がともに現れる．

同時非負性には $d\ge1$ が必要で，
$0\le D_a\le d$ から
\[
 (s-1)d+b-1\le k\le sd+b-1
\]
を得る．$k=b-1$ なら $s=1$，$u=t$ となり，
線形系の第三式は
\[
 t(b-1)T-tbU=b(b-1)>0
\]
となる．しかし $U=S+T+d-1>T$ なので左辺は負である．
$k=b$ なら $(s-1)d\le1$ であり，$s=1$ または $(s,d)=(2,1)$ となる．
前者では第三式 $tb(T-U)=b^2-1>0$ が同じ理由で矛盾する．
後者では $\alpha=b+t$，$U=S+T$ であり，第一式と第三式の線形結合が
\[
 t(1+b^2+tb)S=t+2b-b^3<t
\]
を要求するため，$S\ge1$ に反する．
この境界排除は determinant で割らないので，途中で特異系を落としていない．

以上から $k\ge b+1$ であり，
$D_c\ge0$ と $k\le sd+b-1$ から
$\alpha-k\ge(t-1)d\ge0$ が従う．
$\bar b=-\mathcal B=s+tb$ と置けば
\[
 \mathcal C-\bar b=(s-1)b+u(k-b)\ge0,\qquad
 \Delta=\bar b(\bar b-1)+(\mathcal C-\bar b)(\mathcal A+\bar b)
\]
により determinant の狭義正値性を得る．

残る非特異領域での一般符号問題は未解決である．
この領域の整数局所模型自体は存在する．例えば
\[
 (s,t,\alpha,\beta,k)=(1,2,10,-2,5),\qquad
 (S,T,U)=(4,5,11)
\]
は $q=(5,11,23,1)$，$D=(-1,1,3,-3)$，
正の weights $(5,109/6,45/22,1)$ を与える．
従って $\alpha\ge k\ge b+1$ と $\Delta>0$ だけで局所模型を捨てず，
Cramer 整除性と全高 $D$ の非負性を引き続き課す必要がある．
\nolinkurl{verify_high_d_boundary_v20.py} は
四十二の多項式恒等式と三つの厳密行列回帰を検算する．
一般命題の根拠は上の恒等式・不等式であり，有限範囲の代入結果ではない．

\paragraph{五つの整数を固定した三高 $D$ の完全復元．}
\label{nhd21:reconstruction}
同時非負性の残余問題では，$x=D_a,y=D_b,z=D_c\ge0$ を直接指定する．
$s,t$ も固定することが必要であり，$D$ だけを固定した有限性とは主張しない．
次の変数は再開後の復元と消去にそのまま使える：
\[
 u=s+t-1,\quad d=x+y+z\ge1,\quad
 k=sd+b-1-x,\quad\alpha=ud+b-1+z,
\]
\[
 h=s-x,\quad j=t-y,\quad\ell=u+z,\quad
 \lambda=ux+sz,\quad\mu=uy+tz,
\]
\[
 F=b-1+\lambda,\quad m=F^2+F+1,\quad
 E=\mu h-\ell-j,\quad J=h+\ell-\lambda j.
\]
正の模型では $h\ge1$，すなわち $D_a\le s-1$ が新しく必要となる．
$j$ の符号は制限しない．$E,J$ は $b$ に依存せず，
\[
 S=h+\frac{FE-J}{m},\quad
 T=j+\frac{E+(F+1)J}{m},\quad
 U=\ell+\frac{(F+1)E+FJ}{m}
\]
が全尺度を一意に復元する．三尺度の整数性は単一条件 $m\mid FE-J$ と同値である．
元の三つの高高 pair 残差はこの尺度で互いに等しくなるが，自動的に零にはならない．
\[
 (-s-tb)S+(s\alpha-t)T-\alpha b=0
\]
を最後まで保持し，全尺度の正値性，canonical 六符号量の正値性，
$\alpha\ge k\ge b+1$ も確認する．これらは元の模型への必要十分判定となり，
正の weights は公開稿の明示式から得られる．

\paragraph{ノルムによる全 $b$ の有限化．}
$v=S-h,w=T-j$ と置くと
\[
 E=(F+1)v+w,\qquad J=-v+Fw,
\]
\[
 N=E^2+EJ+J^2=m(v^2+vw+w^2).
\]
零ノルムは，正尺度の下では
$(s,t,x,y,z)=(3,3,0,0,1),(2,2,1,1,1)$ に限られる．
元の残余式がそれぞれ $b(b+5)=27$，$b(b+9)=14$ を要求するため，両方とも排除される．
従って $N>0$ であり，
\[
 m\mid N,\qquad m\le N,\qquad
 b\le 1-\lambda+
 \left\lfloor\frac{\sqrt{4N-3}-1}{2}\right\rfloor
\]
となる．実装では平方根を \texttt{isqrt} で厳密に評価する．
$s,t,x,y,z$ を固定すれば，任意 cutoff なしに $b\ge2$ を全て調べられる．
この五つの整数そのものへの一様上界は未証明である．

\paragraph{誤った簡約を防ぐ二つの反例．}
$(s,t,x,y,z,b)=(1,1,0,3,0,3)$ では
$F=2,m=7,E=4,J=2,N=28$ である．$m\mid N$ でも
$m\nmid FE-J=6$ なので尺度は整数でない．ノルム整除性だけに置き換えない．
また $(1,3,0,4,0,2)$ では $(S,T,U)=(3,5,11)$ は正整数で，
六符号量も正，$D=(0,4,0,-4)$ であるが，元の三つの高高 pair 残差は全て三である．
この例は模型ではなく，残余式を落とせないことの確認例である．

\paragraph{再利用する実装と次の課題．}
\nolinkurl{verification/verify_high_d_reconstruction_v21.py} の
\texttt{reconstruct\_fixed\_data(s,t,x,y,z)} が全 $b$ の完全復元器である．
通常の実行は二十四の多項式係数恒等式と五つの厳密回帰を確認する．
\texttt{--fixed 3 4 0 1 1} のように固定データを渡せる．
既存の有限族検査は $s,t\le12,0\le x<s,0\le y,z\le4,d\ge1$ の
$23256$ 組について $b\ge2$ を全範囲で尽くし，模型は零だった．
結果と実行当時の source は \nolinkurl{research/} に保存した．
これは五変数に制限を付けた結果であり，一般不存在ではない．
次は単一合同式と元の残余式を同時に使い，五変数側の無上界障害または反例を探す．


\paragraph{整数性を外すと，三つの高 $D$ が全て狭義正になる．}
次は\emph{有理数による代数的緩和}の反例であり，整数 branch denominator を持つ
局所 chart の反例ではない．頂点順を $(a,b,c,0)$ とし，
\[
 q=\left(\frac{79}{25},\frac{128}{25},\frac{274}{25},1\right),\qquad
 P=\begin{pmatrix}
 0&179/25&-231/25&-27/25\\
 53/25&0&-259/50&-103/50\\
 -91/25&-266/25&0&83/25\\
 -2&-2&3&0
 \end{pmatrix}
\]
と取る．全 row sums と pair equations を満たし，
\[
 w=(1,7/3,13/4,1)
\]
について全ての $w_j-w_ip_{ij}/q_i$ は狭義正である．
実際，正の辺は $a\leftrightarrow b$ と $c\leftrightarrow0$ のみで，
対応する四つの差は
\[
 \frac{16}{237},\quad\frac{13}{384},\quad
 \frac{17}{1096},\quad\frac14
\]
である．一方，
\[
 D=\left(\frac9{25},\frac9{25},\frac{23}{50},-\frac{59}{50}\right).
\]
この例でも low row と Schur 補行列
\[
 B=\begin{pmatrix}1&5&-6\\-2&1&1\\3&-4&1\end{pmatrix}
\]
は整数である．従って，$B$ と low row の整数性まで仮定しても，
$q$ と逆向き numerators の整数性を外すと一般符号原理は成立しない．
正の weights と主小行列の不等式だけによる証明は不可能であり，
残る整除性を本質的に用いる必要がある．
この反例と上記恒等式は \texttt{verify\_three\_high\_schur\_v12.py} で有理数演算により再確認できる．

\paragraph{$\beta=-1$ を持つ無限族と，一つの slice の完全な記述．}
\label{nsch18:beta-minus-one}
上の完全約数表示の $t=1$ の場合には，さらに簡単な閉じた表示がある．
一般のmatching型に $\beta\leq-2$ を仮定してはいけないことも，この族から分かる．
実際，$t=1$，$\beta=-1$，$k=0$ と固定した完全マッチング型は，
次の二つの正整数 $d,e$，$de\geq2$ でちょうど記述される：
\begin{align*}
 s&=de,&\alpha&=d(e+1),\\
 S&=d(e+1)^2(de-1),&
 T&=de(e+1),&
 U&=(e+1)(de-1).
\end{align*}
分母と numerators は
\[
 (q_a,q_b,q_c,q_0)=(1+sS,1+T,1+sU,1),
 \qquad
 P=\begin{pmatrix}
 0&S+\alpha&-(sS+\alpha+1)&-S\\
 sT-1&0&-sT&-T\\
 -sU&-(U+1)&0&U\\
 -s&-1&s&0
 \end{pmatrix}
\]
である．三つの高分母は全て奇数で，一より大きい．
$H=S+\alpha$，$K=sT-1$ と書くと，正の weights は例えば
\[
 w=\left(q_a,\ H+\frac1{2K},\ s+\frac1{2U},\ 1\right)
\]
で与えられる．正の四つの directed edges に対応する basin gaps は
\[
 \frac1{2K},\quad\frac1{2q_b},\quad\frac1{2U},\quad\frac1{2q_c}
\]
であり，残る edges は負の numerators を持つので正値性は自明である．
また
\[
 D_0=S+T-U-1=(\alpha-1)U+T-1>0,
 \qquad
 (D_a,D_b,D_c)=(-sD_0,-D_0-\alpha,sD_0+\alpha).
\]
従って，これは三つの高係数が全て正となる反例ではない．

完全性の確認も短い．$t=1,\beta=-1,k=0$ の線形系では
\[
 \Delta=s(s-\alpha),\qquad
 S=\frac{\alpha^2(s-1)}{\alpha-s},\quad
 T=\frac{\alpha s}{\alpha-s},\quad
 U=\frac{\alpha(s-1)}{\alpha-s}.
\]
$s\geq2$ であり，$\alpha=s$ の場合には $S$ の Cramer numerator
$-\alpha^2s(s-1)$ が非零なので系は不整合である．
$S>0$ には $d=\alpha-s>0$ が必要である．
$T,U$ の整数性は $d\mid s^2$ および $d\mid s(s-1)$ と同値なので，$d\mid s$ が従う．
逆に $d\mid s$ なら三つ全てが整数となる．$s=de$ と書けば表示した族を得る．
pair equations と row sums はそのまま代入して確認できる．

小さい具体例として $d=1,e=2$ を取ると
\[
 q=(19,7,7,1),\qquad
 P=\begin{pmatrix}
 0&12&-22&-9\\
 11&0&-12&-6\\
 -6&-4&0&3\\
 -2&-1&2&0
 \end{pmatrix},\qquad
 D=(-22,-14,25,11),
\]
かつ $w=(19,265/22,13/6,1)$ となる．
これは既存の有限 catalog にも含まれる模型であり，この slice 全体の無限パラメータ表示と，残る整除性の特徴づけを再利用できる．
これらの恒等式と具体例は \texttt{verify\_three\_high\_bounds\_v13.py} で厳密演算により検算する．

\subsection{次の接着計算に使う二つの局所障害}
\label{ng16:local-obstructions}

二価 core $i$ の二つの高近傍を $j,k$ とする．$q_{jk}=1$ なら，
core $i$ の二高分母標準形は必ず正の枝である．従って
\[
 S>r,\qquad rT>(r+1)S,\qquad L=T-S+1\geq3.
\]
この結論に対向 $D$ の相殺は不要であり，core $j,k$ の高分母次数を制限しない．
負の枝 $S=T=w<r$ を用いるときは，その core を含む高分母三角形が必要である．
一般の三高分母中心に二高分母の真約数補題を適用する代わりに，
中心側では signed divisors と高辺の合同式を直接比較する．
負の枝の $H=r-w$，$c=s-x$，$e=z-y$ に対して
\[
 Ece=e-c+1,\qquad E=1-\frac{w^2}{H(H+1)},
 \qquad -1<E<\tfrac12,\quad E\ne0
\]
となり，中心間が低辺なら reverse 和が $E-1$ となることが主な障害である．
例外 $e=1$ も node 数 $2+rw/H>6$ により排除される．

全低分母 core $*$ がある場合には，残る四頂点を $I$ とし
\[
 x_{ij}=p^i_{j*},\qquad r_{ij}=p^i_{*j}
\]
と置く．四辺 $i*$ の $D$ 相殺，局所の行和，および高辺の整数 clutching から
\[
 \sum_{i\ne j\in I}x_{ij}=\sum_{i\ne j\in I}r_{ij}=-4,
 \qquad
 \sum_{\substack{i<j\in I\\q_{ij}=1}}
 (x_{ij}+x_{ji}-r_{ij}-r_{ji})=0
\]
を得る．全低 core の自由な entries を固定しなくても使える総和条件である．
diamond では和に低辺が一つだけ残り，二価 core の低低 entries が
反対向き・同じ $L\in\{1,3,4\}$ にそろう．三つの場合は全て排除済みである．
$K_4$ 型ではこの和は空になるが，四本の $D$ 相殺を頂点別に保持すると，
後述する通常辺上の流れの式を得る．


\subsection{共通低 puncture における通常接着の上界}
\label{ng16:ordinary-bound}

$q_{ij}>1$，$q_{ik}=q_{jk}=1$ とし，高辺 $ij$ の puncture $k$ が通常の反対符号接着なら
\[
 p^i_{jk}=-p^j_{ik},\qquad p^i_{kj}=-p^j_{ki}=u,\qquad1\le |u|\le4.
\]
core $k$ の低低 entries を $a=p^k_{ij},b=p^k_{ji}$ とすれば $ab=0$ であり，
二つの node 数 $n,m\in\{0,\ldots,5\}$ に対して
\[
 a=1-n-u,\qquad b=1-m+u,\qquad
 0=-\{u-(1-n)\}\{u-(m-1)\}
\]
となるためである．高分母次数と対向 $D$ の相殺には依存しない．
node 総数が同じ $N\ge2$ の模型では，同じ証明で $|u|\le N-1$ を得る．
最大・単位約数の同符号例外をこの上界に含めてはいけない．

\subsection{一高分母の自由度を保持する標準形}
\label{ng16:leaf-data}

一高分母 core は，高 ray $A$ と低 rays $B,C,D$ の適切な標識により
\[
\begin{gathered}
 m=th(h+1),\quad q=m+1,\quad u=th,\quad v=t(h+1),
 \quad h,t\geq1,\\
 P=\begin{pmatrix}
 0&-m&-h-1&h\\
 -1&0&0&0\\
 -u&u-1-c&0&c\\
 v&-v-1-d&d&0
 \end{pmatrix},\qquad c,d\in\ZZ,\quad cd=0
\end{gathered}
\]
と書け，全ての整数 choices に正の weights が存在する．係数は
\[
 (D_A,D_B,D_C,D_D)=(-m-t,\ m+t+1+c+d,\ h-d,\ -h-1-c).
\]
$q=3$ では二つの負の reverse entries が一致し得る．
高 row だけを固定して低低 entries を一つの模型に置き換えず，
この $c,d$ の自由度と必要な低 ray の交換を保持する．

二高分母の低係数は
\[
 D_C=y-x+T-S,\qquad D_D=s-z-1
\]
であり，同時に負とはならない．一高分母の負 reverse を持つ ray $E$ は $B$ または $C$ である．
別の低 ray $J$ から $p_{JE}\leq-1$ を要求される場合，
$E=C$ なら $D_E>0$，残る低 ray から $E$ への entry は零となる．
$E=B$ かつ $D_E\leq0$ なら，残る低 ray からの entry は
$t(h+1)^2-D_E\geq4$ である．
係数相殺と node 条件を組み合わせるときに使える二分である．

\subsection{二葉付き三角形の完全な有限化から再利用する式}
\label{ng16:two-leaf-certificate}

$H=\{01,02,12,03,14\}$ は全排除済みである．再利用するのは，
未知の中心分母を走査せず，葉の自由度を保存したまま有限化する次の手順である．
中央 core $2$ の交差標識 $(A,B,C,D)=(0,1,4,3)$ と正の枝に対して，
共通低 puncture の上界は
\[
 1\le r\le4,\qquad z=1
\]
を同時に要求する．$T\ge5$ により辺 $12$ は正の最大・単位約数の例外となり，
\[
 b=T+1,\quad p^1_{23}=T,\quad p^1_{32}=1,\quad
 p^3_{12}=0,\quad n_{12;3}=n_{13;2}=0
\]
を得る．$r>1$ なら，辺 $02$ には通常接着が強制される．

$T=(r+1)y$，$E=D^2_4$ と置くと，
\begin{align*}
 (r+2)sS-r(rs+1)y&=2r(s-1),\\
 F_r(s)S&=r(rs+1)E+2r(r+2)(s-1),\\
 F_r(s)&=-rs^2+(4r+3)s-r,\\
 E&=(r+2)y-\frac{s+r}{r}S.
\end{align*}
$r=1,\ldots,4$ では $F_r(s)\ne0$ であり，$E\ge-2$ なら $s\le9$ となる．
より一般に，整数の正枝で $z=1$ なら，任意の正整数 $r,s$ でも $F_r(s)\ne0$ である．
実際，$F_r(s)=0$ は $r=3s/(s^2-4s+1)$ を与え，正整数解は $(r,s)=(12,4)$ だけだが，
上の係数式は $E=-12/7$ を要求して整数性に反する．
従って $(r,s,E)$ は $S$，続いて $y,T$ を一意に復元する．
また $E\ge-M$，$M\in\ZZ_{\ge0}$ なら
\[
 E<\left(4+\frac4r-s\right)\frac Sr,\qquad S>r
\]
から $s\le M+7$ を得る．これは他の係数下界が得られた場合にも再利用できる．
$S\ge4$ の下で，通常の $C/D$ 表現を持つ葉 $4$ では，
$-2\le E\le7$ の枝と，$r=1$，$D^4_3=E+\kappa$，
$\kappa\in\{2,5,8,10\}$ の枝に分かれる．
後者でも葉 $3$ の一方の低標識なら $th\le4$，
他方なら $h\le12,t\le6$ が証明される．
小さい $S$ の共鳴族 $(r,S,T)=(2,3,6)$ では，$s$ は正の偶数であり，
$s\ge8$ は他葉に向かう二つの正の $D$ が相殺しないため排除される．

この有限化は三十一中央模型・六十六生成 witness を与え，
全葉を復元すると十中央模型・五十二低三角形だけが残る．
葉 $4$ の ray $2=B$ を含む別の有限化は，
三十三中央模型・六十二 witness から四中央模型・七組を与える．
この七組は五十二組の部分集合であり，個数を加えてはいけない．
分母三の同じ標識付き行列は，一度だけ数える．
$B$ 表現の例外を省くことも，低低自由度 $c,d$ を固定することも許されない．

\paragraph{既知三行だけによる中心の延長障害．}
三つの隣接 cores の行列が既知なら，中心 $i$ の row $j$ は
\[
 p^i_{jk}=q_{ij}(1-n_{ij;k})-p^j_{ik}
\]
で復元される．非負 node 数の三分割は二十一通りなので，既知三行は全 $21^3$ 通りである．
三つの既知 rays の pair equations は
\[
 p^i_{jk}p^i_{kj}=q_{ij}q_{ik}-1
\]
である．五十二組のうち五十一組は中心 $0$，残る一組は中心 $1$ において，
この三行条件だけで候補が零となる．
未知の中心間分母，第四行，中心での $D$ 相殺，正の weights の判定は，
この最後の矛盾に必要ない．
この手順は，他の残存 graph で複数の隣接行が既知となる場合にも使える．

\subsection{\texorpdfstring{$K_4$}{K4} の低 hub と通常辺上の流れ}
\label{ng16:k4-flow}

全低 core を $4$ とし，$I=\{0,1,2,3\}$ の高 $K_4$ で
\[
 R_{ij}=p^i_{4j},\quad X_{ij}=p^i_{j4},\quad V_{ij}=p^4_{ij},
 \quad N_{ij}=1-R_{ij}-V_{ij}
\]
と置く．各行の node 予算は $\sum_jN_{ij}=5$ である．
通常辺では $R_{ji}=-R_{ij}$，$1\le|R_{ij}|\le4$，
例外辺では $R$ の対が $\pm\{1,q-1\}$ となる．
hub に接する $D$ 相殺を頂点別に保持すると
\[
 \sum_{j\ne i}(X_{ij}-R_{ji})
 =\sum_{j\ne i}N_{ji}-5
\]
を得る．例外辺の寄与は零であり，通常辺の寄与は
\[
 \operatorname{sgn}(R_{ij})W_{ij},\qquad
 W_{ij}=\frac{q_{ij}-1}{|R_{ij}|}+|R_{ij}|
\]
である．proper cut の発散の絶対値は高々十なので，通常 graph の橋では
$W\le10$，従って $q\le25$ となる．

例外辺数は二・四・六のいずれかであるが，三場合は全て排除された．
二例外は，負の大きい reverse を持つ頂点から見た正例外の位置による三場合で，
hub の対向 $D$ 式と正 flow の下界が矛盾する．三高分母の一般分類は不要である．

四例外の再利用できる有限化は，例外の大きい reverse $m=q-1$ がその大きい端の
一つの行だけに現れることを用いる．異符号二項の差を固定した無限方向を
\[
 m_1=h+\max(\delta,0),\qquad m_2=h+\max(-\delta,0),\qquad h\in2\ZZ_{>0}
\]
と書き，$h=2,4$ と全 $h\ge6$ に分ける．最後の半直線では大きい端の
$V=1-R-N$ が必ず非零なので，pair 条件は単位端の node を正例外で零，負例外で二に固定する．
この後の row，node，hub pair，低辺の $D$ 条件は $h$ に依存しない．
従って $h=6$ は無限型全体の厳密な代表であり，分母cutoffではない．
通常二辺が端点非共有の場合の百十二型では候補零，共有の場合の二百四型では候補六となり，
残る無限型は零である．全六候補に
\[
 (q_1,q_2,q_3)=(5,3,3),\quad r=(-4,1,2),\quad x=(-1,2,1)
\]
のcoreが現れ，残る局所条件は $6A^2+17A+14=0$，判別式 $-47$ となる．
これは実数の高高entriesも許さない．

有限化と独立監査には，reverse-first と node-first の二方向がある．
後者は全 $21^4$ node 行列から通常forestのflowを決め，例外分母の行方程式を別に解く．
両方法は六候補の全行列で一致した．次のgraphでは，一方の表現で無限方向が見えにくい場合に
他方から復元し，同じ有限候補へ至るかを調べる価値がある．
全低coreを持つsupportはこれで全分母一に限られるが，一高分母hubへ移すと
pair productsとrow予算が変わるため，補正を導かずに同じ式を適用してはいけない．

全分母一のatlasや一高分母hubとの比較のため，全低 hub 自身の自由度の表示も残す．
$V_{ij}V_{ji}=0$，行和 $-1$ から，各行に負の非零 entry が少なくとも一つある．
非零 support は四，五，六辺であり，出次数列はそれぞれ
$(1,1,1,1)$，$(2,1,1,1)$，および $(3,1,1,1)$ または $(2,2,1,1)$ である．
出次数一の entry は $-1$，出次数二の行は $(t,-1-t)$ なので，
有限個の有向 support と高々二整数の族として扱える．
特定の全 $-1$ 模型に固定せず，正の weights が必要な段階で cycle 条件を課す．
\subsection{部分的な低辺条件だけでは完成しない二つの模型}
\label{ng16:partial-patches}

次の例は，部分的な必要条件の非空性を示す資料である．大域 atlas の存在例ではない．
$K_{2,3}$ の三つの二価 core を cyclic に標識し，canonical 高 ray $A$ が
一方の中心を向く場合を $+$，他方を向く場合を $-$ とする．
符号 $(+,+,-)$ に対し，canonical データを
\[
 (r,s,z,S,T)=(4,4,5,2,2),\quad(5,10,2,3,3),\quad(5,10,2,3,3)
\]
と取る．三低辺の node triples は $(0,3,2),(2,2,1),(0,2,3)$ だが，
同じ順の対向低係数和は $7,-10,5$ であり，相殺しない．
中心二 core も与えていない．本版の $K_{2,3}$ 排除は，中心間の低辺と高辺の clutching を加えて成立する．
具体的な全行列・係数・node の再計算は，保存した v14 の node 検算器に含まれる．

diamond の二価 cores $2,3$ と全低 core $4$ についても，
三つの相互低辺の node 条件と $D$ 相殺を全て満たす部分模型がある．
行列の添字順をそれぞれ $(0,1,3,4),(0,1,2,4),(0,1,2,3)$ として
\[
 P_2=\begin{pmatrix}0&-6&1&-2\\-8&0&-2&3\\6&-3&0&-4\\-3&2&0&0\end{pmatrix},\quad
 P_3=\begin{pmatrix}0&-14&-2&5\\-7&0&2&-4\\-5&4&0&0\\2&-2&-1&0\end{pmatrix},
\]
\[
 P_4=\begin{pmatrix}0&0&-1&0\\2&0&0&-3\\0&-1&0&0\\-3&0&2&0\end{pmatrix},
 \quad q_2=(7,7,1,1),\quad q_3=(11,9,1,1),\quad q_4=(1,1,1,1).
\]
係数は順に
\[
 (-2,0,0,2),\qquad(-1,3,0,-2),\qquad(0,0,-2,2)
\]
であり，低辺 $23,24,34$ の node triples はそれぞれ
$(0,0,5)$，$(4,0,1)$，$(2,3,0)$ である．
全てに正の weights がある．しかし全低 hub の残差
\[
 p^2_{34}+p^3_{24}-p^2_{43}-p^3_{42}=-3
\]
は零でないので，二つの中心 core を加えた大域完成は不可能である．
行和・pair products・正 weights・係数・nodes の厳密検算は
\nolinkurl{diamond_obstruction_verification_v15.json} に保存した．
部分模型の存在を保ちながら，追加する大域条件のどこで矛盾するかを区別するための例である．


\subsection{house 型 \texorpdfstring{$G_{5,5}$}{G5,5} の全排除から再利用するもの}
\label{nh18:house}

$H=\{01,02,03,12,14,34\}$ は，分母上限なしで排除した．
公開算術稿の定理 \texttt{hs18:no-house} が証明を与える．
二価 core $2$ の負の枝は途中で残し，外側 $3,4$ だけを
「高隣接頂点間が低辺」という既知の補題で正の枝へ制限した．
三価 cores $0,1$ の分類は不要であり，$D$ 相殺を使うのは低辺 $23,24$ だけである．
この型の排除と，現在の五型の存在・不存在の問題は区別する．

\paragraph{共通低 ray の $C/D$ 型による分類．}
正の二高分母標準形の低 rays を $C,D$ とすると，高 ray からの
outgoing / reverse entries は
\[
 (C,A):(x,r),\quad(C,B):(-y,-r-1),\quad
 (D,A):(-s,-S),\quad(D,B):(z,T)
\]
である．正の枝では $x\ge2,S\ge2,T\ge5,L=T-S+1\ge3$ であり，
$y,s,z$ はそれぞれの高分母から一を引いた数の真約数である．
共通低 puncture の signed-divisor 接着と reverse 上界四を組み合わせると，
両 $D$ の型は不可能であり，$C/D$ 混在型では
\[
 \begin{array}{c|c}
 \text{通常}&\text{両 shared rays}=A,\quad r_C=S_D\le4\\
 \text{正例外}&\text{shared ray}=A_C,B_D,\quad r_C=z_D=1
 \end{array}
\]
だけが残る．これは別の graph の隣接二価 cores にも再利用できる．
canonical $D$ row の低低 entry は $-L$ であり，零と取り違えない．

\paragraph{最終的に現れる完全な必要族．}
house 型で中央 core $2$ の高 rays を $A=0,B=1$ とすると，node 条件は唯一の族
\[
 \begin{array}{c|c|ccccc}
 \text{core}&(A,B,C,D)&r&s&z&S&T\\\hline
 2&(0,1,4,3)&1&2t-1&t&2&6\\
 3&(0,4,1,2)&1&1&1&2&6
 \end{array}
\]
へ帰着する．$t\ge1$ は任意の整数であり，探索上限ではない．
$D^2_3=t-2,D^3_2=-1$ なので低辺 $23$ の相殺で $t=3$ となる．
core $4$ は $(A,B,C,D)=(3,1,2,0)$，$a=7,r=1$ で
$D_C=-3$ を要求されるが，この固定データの全七模型は
\[
 D_C\in\{4,9,15,18,19,23,40\}
\]
を満たすので不可能である．
固定 $b=7$ にある共鳴族を捨てたのではなく，この場所の共有 ray は
canonical $A$ と先に強制されており，その共鳴族は適用範囲が異なる．

\paragraph{検算と今後の使用範囲．}
\nolinkurl{verify_house_obstruction_v18.py} は，共通低 ray の全十六型，
node の不等式，物理標識付き node 式と無限族の二十五個の多項式恒等式，
固定 $a=7,r=1$ の全約数を確認する．
最後の七模型は，証明で導いた上界を使う別の original-matrix 復元でも一致する．
検算結果は \nolinkurl{house_support_obstruction_verification_v18.json} に保存する．
これは一般の node 不等式による場合分けと完全な有限計算を組み合わせた証明であり，
数値上限以下の分母を試した経験的排除ではない．
house 型自体の探索は終了し，共通低 ray の型分類と低辺係数相殺を残る型へ移す．

% Reusable results only. Detailed proof: gten18:no-support in arithmetic paper.
\subsection{\texorpdfstring{$G_{5,10}$}{G5,10} の排除から再利用する方法}
\label{ngten18:reusable}

$G_{5,10}=\{01,02,03,12,13,24,34\}$ は，分母の上限なしに排除された．
算術稿の定理 \texttt{gten18:no-support} を参照する．
核心は，三高分母 core の正符号が完全 matching となる定理を，
二価 core の二つの低辺の node 予算と組み合わせることである．
対象は指定した五次の算術的接着模型であり，任意の fibration の分類ではない．

頂点 $4$ の二高分母標準形を $(A,B,C,D)=(2,3,0,1)$ とし，
$a=sS+1,b=zT+1,x=sS/r,y=zT/(r+1)$，$\delta=T-S$ とする．
$q_{23}=1$ なので正の枝であり，$\delta\ge2$ である．
低辺 $04,14$ の node 数をそれぞれ $(n_1,n_2,n_3)$，$(m_0,m_2,m_3)$ とすると，
隣接する三高分母 cores の低 rows が
\[
 R^0=(1-n_1,1-r-n_2,r+2-n_3),\quad
 R^1=(\delta+2-m_0,S+1-m_2,1-T-m_3)
\]
に固定される．各 row には正成分がちょうど一つあり，各 node 組の和は五である．
共通低 puncture $(01;4)$ の三つの接着型は次のように処理する．

\begin{enumerate}
\item 通常型では，有限な node 予算と matching 性だけで
$(r,S,T)=(2,3,5)$ が強制される．整数 pair equation は
$(u-2)(v+1)=-3$ に帰着し，唯一の中央模型は
$(r,s,z,S,T)=(2,2,6,3,5)$，$(a,b)=(7,31)$ である．
しかし反対側の core では，$30$ の signed divisor が
$27\pmod{31}$ であることを要求される．その候補 $27,-4$ はどちらも不適格である．
この部分は $D$ 相殺を使わない．

\item 正例外型では $R^0=(1,-1,-1)$ となる．二つの低辺の $D$ 相殺から
$x-y>0$ と $s-z\le1$ が従う．正の枝の不等式は $s-z\ge1$ を与えるので
$s=z+1$ である．反対側の低 column の二つの負成分がともに $-1$ となり，
$q_{12}+q_{13}=q_{01}+2$ を要求するが，奇分母の parity と矛盾する．

\item 負例外型では $q_{01}\in\{3,5\}$，$\delta=2$ となる．
$h=S-r$，$\alpha=sh/r$，$\beta=z(h+1)/(r+1)$ とすると
\[
 \left(\frac{r^2}{h}-\frac{(r+1)^2}{h+1}\right)\alpha\beta
   =\alpha-\beta+1.
\]
第三角の signed-divisor clutching により，$\alpha,\beta>1$ なら括弧が整数で，
その絶対値は一未満となる．零の場合は $h=r^2/(2r+1)$ となり不可能である．
$\alpha=1$ の例外も保持して処理すると，必ず $\beta=1$ となる．
残る equation は $(h+1)z(z-2)=h-1$ であり，中央模型全体が
\[
 r=3,\quad S=4,\quad T=6,\quad z=2,\quad s=3\alpha,\quad
 a=12\alpha+1,\quad b=13
\]
という無限共鳴族へ帰着する．これを探索上限で切らない．
低辺 $14$ の $D$ と奇分母から $\alpha$ は正の偶数となり，
$q_{12}\ge15\alpha+1$ を得る．低辺 $04$ から $q_{02}\ge8\alpha-7$ を得る．
辺 $24$ の node と matching 性は二つの正整数 $M,N$ に
\[
 M+N=24\alpha+1,\quad M>3q_{02},\quad N>q_{12}
\]
を要求するため，$M+N>39\alpha-20\ge24\alpha+10$ と矛盾する．
\end{enumerate}

再利用すべき点は，低辺の有限な node 予算から隣接三高 core の低 row を固定すること，
matching によってその符号を制限すること，そして第三角の signed divisors で
共鳴族を特定してから $D$ と node の不等式で全族を排除することである．
検算の入口は \nolinkurl{verify_g5_10_obstruction_v18.py} であり，
係数ごとの恒等式二十四本と，証明由来の有限な node の場合分けを保持する．
有界分母探索を一般証明の代わりにしていない．

% Reusable v18 notes; no exploration diary.
\subsection{一高分母 hub の有限型復元と \texorpdfstring{$K_4$}{K4} 一葉型の解決}
\label{nl18:leaf-hub}

高分母 support が $K_4$ 上の $0,1,2,3$ に高辺 $04$ を加えた型 $G_{5,12}$ は排除された．
入力は一高分母の完全標準形と，三高分母 core の正符号 graph が完全 matching となる定理である．
全低 hub の $K_4$ 障害を部分 graph へ拡張したものではない．
$G_{5,7},G_{5,9}$ は別の配置解析で排除された．
この $K_4$ 一葉型の証明を，隣接 core の次数が異なる二型へ無条件に転用したわけではない．

\paragraph{再利用する補正公式．}
任意の hub $4$ について $\ell_i=q_{i4}$ とし，
$R_{ij}=p^i_{4j}$，$X_{ij}=p^i_{j4}$，$V_{ij}=p^4_{ij}$，$N_{ij}=n_{i4;j}$ と置けば，
\[
 R_{ij}=\ell_i(1-N_{ij})-V_{ij},\qquad
 R_{ij}X_{ij}=\ell_iq_{ij}-1,\qquad
 \sum_{j\ne i}N_{ij}=5.
\]
対向 $D$ 相殺は
\[
 \sum_{j\ne i}(X_{ij}-R_{ji})
 =\sum_{j\ne i}\ell_jN_{ji}-\sum_j\ell_j-\ell_i
\]
となる．一高 hub では $\ell=(Q,1,1,1)$ である．
高 ray を含む hub pair は $Q-1$ であり，零積へ置き換えない．

\paragraph{葉分母の上界を導く仕組み．}
標準形を $(A,B,C,D)=(0,1,2,3)$ とすると
$u=th$，$v=t(h+1)$，$Q=th(h+1)+1$ である．
特別な $B$ ray の低 row は $R_{10}=2-N_{10}$，$R_{12},R_{13}\le1$ となる．
一正二負条件と clutching から $R_{10}=1,2$ のみが残り，
$q_{01}=2Q-1,Q,2Q+1$ の三族と，$Q=3,5$ における任意奇数 $q_{01}$ の退化族を得る．
$2Q-1$ 族は高 hub 辺の $D$ と他の高辺の node 条件により $Q=3$ に限られる．
$Q>21$ なら $u\ge5,v\ge6$ である．$d=0$ は $R_{31}\ge3$ と
$R_{13}\in\{-1,-2\}$ の反対符号接着に反する．$c=0$ では
core $2$ の唯一の正 entry が $R_{20}$ となり，$|R_{23}|\le4$ である一方，
core $3$ の row sum と一正二負条件は $R_{32}\ge5$ を要求する．
従って $Q\le21$ である．

\paragraph{自由整数と他の分母を切らない復元．}
$Q\le21$ の $(h,t)$ は十五個，四つの node 行はそれぞれ二十一通りである．
$c=0,d=z$ または $d=0,c=z$ とし，$z$ は全整数にわたって残す．
辺 $0i$ では $R=R_{0i},r=R_{i0}$ が $z$ に依存せず，
\[
 (Qr+R-\kappa Rr)q_{0i}=R+r,\qquad -4\le\kappa\le1
\]
から固定分母または全無限族 $q_{0i}=1+\operatorname{lcm}(2,|r|)a$ を得る．
低三角形内では固定 reverse $r$ の相手は $-r$，条件付きの $\operatorname{sgn}(r)$，
または $|r|=1$ における同符号の全正偶数である．
従って $z$ の許容集合は点と偶奇付き半直線の有限和となり，正確に交差できる．

固定 node 配置 $345$，partner 型の組 $1192$ のうち符号条件を満たす $579$ 組は，
全て空交差である．これらは手順内の枝を含む型の数で，charts の同型類数ではない．
\nolinkurl{verify_k4_leaf_obstruction_v18.py} と
\nolinkurl{k4_leaf_obstruction_verification_v18.json} に全 $579$ 個の空交差証明書を保持する．
$Q=3$ における標識重複を残し，$c=d=0$ の二重計数も不存在判定のため許している．
この方法では他の高分母や自由な低低係数に探索上限を置かない．

\subsection{二つの葉付き支持の排除から再利用する方法}
\label{nl19:section}

$G_{5,7}$ と $G_{5,9}$ は，一高完全標準形，三高 matching，
二高の正負二分，共通低 puncture の signed-divisor 分類，node 非負性と対向 $D$ 相殺で排除された．
標準 $K_5$ graph，四直線 cones，marked punctures を保つ五次の整数接着問題に関する結論である．
葉を $4$，その高隣接を $0$ とし，葉の標準 rays を $(0,B,C,D)$ とする．
$Q=th(h+1)+1$，$u=th$，$v=t(h+1)$，$cd=0$ の全整数自由度を保持する．

$G_{5,7}$ では $0$ は三高で，その唯一の低 ray の位置 $m$ が $B,C,D$ の全配置を尽くす．
可能な配置は
\[
 (B,N_{0B}=2,Q\le5),\quad(C,N_{0C}=1,u\le4),\quad
 (C,N_{0C}=2,h=t=1),\quad(D,N_{0D}=1,v\le4)
\]
に限られ，従って $Q\le21$ である．
$G_{5,9}$ では $0$ は四高で，残る三高 core の位置 $J=B,C,D$ を全て残す．
$B$ の特別な reverse は $J=B$ なら $1,2$，他なら $1,2,-2$ だけとなる．
$-2$ の枝は $Q=3$ に限られ，三配置の低低 entry と matching の比較から $Q\le13$ を得る．
いずれの上界も，当該 support の仮定を用いて証明した結果であり，数値探索の上限ではない．

再利用する葉接着の式は
\[
 \ell=(Q,1,1,1),\quad R_{ij}=p^i_{4j},\quad X_{ij}=p^i_{j4},\quad V=P^4,
\]
\[
 R_{ij}=\ell_i(1-N_{ij})-V_{ij},\quad\sum_{j\ne i}N_{ij}=5,\quad
 R_{ij}X_{ij}=\ell_iq_{ij}-1,
\]
\[
 \sum_{j\ne i}X_{ij}+\sum_{j\ne i}V_{ji}=-2\ell_i.
\]
高 ray を含む葉の pair product は $Q-1$ であり，零ではない．
各 node 行は二十一通りだが，それだけでは自由整数も他の分母も有限にならない．
固定 reverse の相手は一点または全偶奇半直線で処理し，中央辺では
\[
 (Qr+R-\kappa Rr)q=R+r
\]
を用いる．係数と右辺がともに零なら，
$q=1+\operatorname{lcm}(2,|r|)a$，$a\ge1$ の全族を残す．

$G_{5,7}$ の残余は $160$ partner 型，中央辺の選択まで含めて $212$ 型である．
中心の $D$ 式に必要な情報をまとめた十三行では，自由分母の七行が $q=1,-1,-3$ を要求し，
残る六行は非零の $D$ 残差を持つ．
$G_{5,9}$ では点 $908$ 型と全半直線 $18$ 型が残る．後者は二高 core の canonical な
第二高分母 $b=3$ を要求して矛盾する．低低 pair の逆 entry は別の低辺の node 予算から
$0,-1,-3,-4$ に限られ，点型から $2596$ 個の $D$ 線形系を得る．
$2563$ 個は有理係数で不整合，残る $33$ 個は $3a+3b=-6$，$a,b\ge1$ を要求する．
各矛盾の元の四式による有理線形結合を証明書に保存した．
自由分母が多数残っても，先に全局所行列を完成させず，$D$ の必要条件だけで排除できることを示す．

\paragraph{小さい正 entry と，局所例の使い分け．}
二高 core の二つの低行の正 entries が同時に二以下になることはない．
canonical には $r,T$ であり，正枝は $T>r+1$，負枝の唯一の小候補は
$(s-4)(z+3)=-6$ と整除性に反する．
この補題は，隣接 matching core から node によって小さい上界が伝わる場合に使える．
一方，三高 core の低行 $(-1,-1,1)$ 自体は不可能ではない．例えば
\[
 q=(1,3,31,7),\qquad
 P=\begin{pmatrix}
 0&-1&-1&1\\-2&0&4&-5\\-30&23&0&-24\\6&-4&-9&0
 \end{pmatrix}
\]
は row/pair equations と正値性を満たす．正二閉路の gains は $6/7$ と $92/93$ である．
葉型の排除には隣接 core と node/$D$ を用いることを忘れない．

完全証明は算術稿の \nolinkurl{g5_7_obstruction_v19.tex} と
\nolinkurl{g5_9_obstruction_v19.tex}，型の全証明書は同名の主検算と結果 JSON にある．
型数は canonical branches の重複を残す手順上の個数であり，同型類数ではない．

\subsection{二つの三価頂点を持つ七高辺型の排除と，再利用する有限化}
\label{ng51120:obstruction}

\[
 G_{5,11}=\{01,02,03,04,12,13,24\}
\]
は，正の整数局所 charts，marked clutching，非負 node，対向 $D$ の相殺の下で存在しない．
二価 cores は $3,4$，三価 cores は $1,2$，四価 core は $0$ である．
この排除には分母および scale の探索上限を用いない．
次の型へ再利用するため，有限化の仮定と無限族の保持を記録する．

\paragraph{向き別の reverse 上界．}
二価 core を $i$，他方の二価 core を $j$，
固有の高近傍を $h$，残る三価頂点を $k$ とする．
canonical order $(A,B,C,D)$ と ray $j$ の行を，列 $(0,h,k)$ の順で書くと
\[
\begin{array}{c|c|c}
\text{型}&(A,B,C,D)&\text{ray }j\text{ の行}\\\hline
\mathrm I&(0,h,j,k)&(r,-r-1,0)\\
\mathrm {II}&(h,0,j,k)&(-r-1,r,0)\\
\mathrm {III}&(0,h,k,j)&(-S,T,-L)\\
\mathrm {IV}&(h,0,k,j)&(T,-S,-L)
\end{array}
\]
となる．共通低 puncture の通常 reverse の絶対値は四以下である．
正例外 $\{1,q-1\}$ と負例外 $\{-1,-(q-1)\}$ を捨てずに比較する．
正の単位 entry $x=sS/r=1$ は，正枝では $S>r$，
負枝では pair equation の符号から一般に排除できる．
これにより I 型は $r\le3$，II 型は $r\le4$，IV 型は $S\le4$ となる．

正の III 型は先に除く必要がある．相手の node と三高 matching により，
両二価 cores が $(r,s,z,S,T)=(1,1,1,2,6)$ に固定される．
外側低辺の $D$ 相殺が $q_{01}+q_{12}=4$ を要求し，両分母三以上と矛盾する．
負の III 型は $w=S=T\le4$ であり，他の負枝も $w\le4$ となる．
$w^2/2<H(H+1)<2w^2$ から $H\le5,r\le9$ を得る．

正枝の相手別の議論は，上界を他の graph へ移す際に確認すべき箇所である．
I--II では相手 $(r',S',T')=(1,2,5)$，
I 側は $(r,S,T)=(3,4,6),(3,5,7)$ に限られる．
II--II では $r=r'=1$，$S,S'\le3$ となり，pair equation が $T,T'<14$ を与える．
II--IV は通常接着なら $T\le5$，負例外なら $L=3$，
$r<S<2r$ から $S\le7,T\le9$ となる．
IV の残る相手は II または負の III であり，それぞれ $T\le6,T\le5$ である．
従って全体は，数学的に証明された有限範囲
\[
 1\le r\le9,\qquad 2\le S\le7,\qquad S\le T\le13
\]
に含まれる．この範囲は全候補を含む上界であって，全ての点が接着可能という意味ではない．

\paragraph{全 scale を保持した点・半直線表示．}
$h=S-r,\delta=T-S$ とすると
\[
 Fsz+As+Bz+C=0,
\]
\[
 F=(2r+1)h-r^2(\delta-1),\quad
 A=h(r+1),\quad B=-r(h+\delta-1),\quad C=r(r+1)
\]
である．$F\ne0$ では
$(Fs+B)(Fz+A)=K=AB-FC$ の全 signed divisors を用いる．
$K=0$ の二直線と $F=0$ の一次方程式は，整数合同式を解いて全半直線として残す．
この graph では book と異なり $L\le4$ を仮定できず，$F=0$ は実際に生じる．
無限族を持つ reverse triples は
\[
\begin{array}{c|c|l}
(r,S,T)&\text{場合}&\text{全正整数 scale}\\\hline
(1,2,6)&F=0&s=2z-1,\ z\ge1\\
(2,6,12)&F=0&s=3a-2,\ z=2a-1,\ a\ge1\\
(2,3,6)&K=0&s=2a,\ z=1,\ a\ge1\\
(3,4,6)&K=0&s=3a,\ z=2,\ a\ge1
\end{array}
\]
である．後二つだけを共鳴として残す手続きでは不完全になる．
整除性と奇数分母条件による剰余類分割の後，一つの二価 core あたり
$90$ 点型と $78$ 全半直線型，計 $168$ 表示を得る．
向き I,II,III,IV の表示数はそれぞれ $57,61,8,42$ である．
これは幾何的模型の個数ではなく，無限族を含む完全なパラメータ表示数である．

\paragraph{最後の低辺と，独立監査の役割．}
低辺 $34$ の node を通る型対は $3592$ 個である．
外側低辺 $14,23$ で各二十一 node 分配を使い，三高の低 row の一正二負，
共通低 puncture の通常・例外接着を課すと $1395$ 配置が残る．
例外接着の要求分母と affine 分母の一致から，
定数不一致 $759$，parameter 非整数 $550$，計 $1309$ 配置を除く．
残る $86$ 配置では低辺 $34$ の残差 $D^3_4+D^4_3$ が
$72$ 配置で常に奇数，$14$ 配置で非零定数となる．
定数の内訳は $8$ が四配置，$28$ が一配置，$-4$ が九配置である．
自由 parameter の上限を設けず，全ての非負整数に対して相殺を排除する．

主検算 \nolinkurl{verify_g5_11_obstruction_v20.py} は
全 affine 行列，node 行，要求分母，残差を証明書に保存する．
独立した \nolinkurl{audit_g5_11_certificates_v20.py} は主生成器を import せず，
canonical row/pair 式と元の node 定義から全 $1395$ 配置を再構成する．
さらに整数性の primitive step を用いる別の生成法で
$90$ 点と $11$ primitive 半直線を得て，
主検算の $78$ 剰余類半直線がその全体を覆うことを有限周期ごとに照合する．
この監査は半直線の最初の有限点だけを照合するものではない．
無上界の結論は解析的 reverse 上界と全整数族の処理の組合せに基づく．

% Research-note module, v21. Keep actionable material and exact scope only.
\subsection{\texorpdfstring{$G_{5,15}$}{G5,15} の継続に使う還元}
\label{ng51521:section}

\subsubsection{今回確定した境界}

低辺を $23,24$ とする標識では，core $2$ が二高，$3,4$ が三高，$0,1$ が四高である．
公開稿 \nolinkurl{g5_15_reduction_v21.tex} の定理 \texttt{g51521:reduction} は，
共通低 puncture $2$ を持つ共有高辺 $34$ の reverse 分岐を完全に分ける．
正例外型は無上界で排除された．通常型は中心二価 core の四点と十四の低 row 配置に限られる．
負例外型は $q_{34}=3,5$ に限られ，node 配分まで明示した二つの template が残る．
特に $q_{34}\ge7$ では，二価 core の二高分母が canonical 順で
$(7,31),(9,11),(31,7),(19,43)$ のいずれかになる．
支持 $G_{5,15}$ 自体は未排除であり，支持型の個数を減らしてはならない．

公開稿は全ての通常型中心，十四配置，正例外の八矛盾，負例外 template を保持する．
これらの表を研究ノートで重複させる必要はない．
検証器 \nolinkurl{verify_g5_15_reduction_v21.py} は旧版を import せず，
全 signed divisors と全 affine 半直線を処理する．
証明ファイルの SHA-256 と検証器自身の SHA-256 を通常版 JSON に記録する．

\subsubsection{通常型で次に解くべき五つの局所問題}

通常型では，共有高 ray への低 row 成分が正である三高 core を先に復元できる．
外部二標識の交換を除くと，低 row $R$ と低 $D$ の組は次の五つだけである．
\begin{center}\small
\begin{tabular}{c|c|r}
中心側の配置&$R$&低 $D$\\\hline
$M_1$, $|u|=4$&$(-1,-4,4)$&5\\
$M_1$, $u=1$&$(-1,-1,1)$&$-9$\\
$M_2$, $|u|=2$&$(-1,-2,2)$&2\\
$M_3$, $|u|=2$&$(-1,-2,2)$&$-7$\\
$M_4$, $|u|=2$&$(-1,-2,2)$&9
\end{tabular}
\end{center}
従って次の有望な目標は，これらの固定 $(R,D_0)$ に対する全正整数模型の解析である．
単に $q_{34}$ の探索上限を増やすより，この五問題の無上界分類または真に無限な族の構成を優先する．
もし三つ目と五つ目だけに有限個の模型が残ると証明できれば，
既存の固定 $q_{34}$ lifting 手順で通常型全体を解ける可能性がある．
有限範囲の疎な観測はその有限性の証明ではない．

計算補助として，$R=(-1,-m,m)$，低 column $(-S,-T,U)$ を取り，
$D_0=S+T-U-1=d$ を課す．高 rays の順を $a,b,c$ とし
\[
 q_a=S+1,\quad q_b=mT+1,\quad q_c=mU+1,
 \qquad U=S+T-d-1
\]
と置く．$H=p_{ab}>0$ なら $p_{ac}=-(H+1)$ である．
\[
 A=q_aq_b-1,\quad B=(m-1)T+1,\quad
 C=q_aq_c-1,\quad D=q_bq_c-1,\quad E=(m+1)U+1
\]
を用いると，残る高 row sum は
\begin{equation}\label{ng51521:simplified-quadratic}
 (D-EB)H^2+\{CB+D-E(A+B)\}H+A(C-E)=0
\end{equation}
に等しい．これは
\[
 \frac C{H+1}+\frac{DH}{A+BH}=E
\]
を払分母した式である．二次係数は
\[
 D-EB=(T-1)(U+1)
\]
と簡単になる．$T=1$ の境界は一次式に退化するため，この式だけで全域を genuine quadratic として扱ってはならない．
$H\mid A$，$H+1\mid C$，$A/H+B\mid D$，全符号，奇分母，および元の row/pair equations も必要である．
この単純化は今後の約数分解や整数不等式の出発点になる．

\subsubsection{固定共有分母の完全 lifting と有限探索の区別}

公開稿の命題 \texttt{g51521:fixed-Q-finite} によって，通常型の十四配置の各々と
任意の固定した奇数 $Q=q_{34}\ge3$ に対し，cores $2,3,4$ の全相互接着を有限に復元できる．
正 reverse 側で $m=|u|$，$U=(Q-1)/m$ とすると，低 $D$ が残る二 scale の和を固定する．
次に辺 $34$ の五つの node 分配を列挙し，反対側の共有高 row を決める．
反対側の最後の scale は共有高 $D$ の一次式で決まる．
一次式が共鳴しても，残る pair equation の二次係数
\[
 \frac{(m-1)^2U+m-2}{\{(m-1)U+1\}^2}>0
\]
が非零なので，高々二候補しか生じない．
この有限性は分母 $Q$ を固定した上での完全性であり，全 $Q$ に共通する上界ではない．

\nolinkurl{ordinary_fixed_q_lifting_v21.py} はこの手順を標準ライブラリのみで実装し，
各正局所模型に正の有理 weights の witness を付ける．例えば
\begin{verbatim}
python ordinary_fixed_q_lifting_v21.py --q 13 --output fixed_q13.json
\end{verbatim}
で，$q_{34}=13$ の十四配置を他の scale cutoff なしで全復元する．
出力に含まれるのは三 cores と三つの相互辺であり，四高 cores $0,1$ は含まれない．

診断用に
\begin{verbatim}
python ordinary_fixed_q_lifting_v21.py --scan-q 1001 --output q1001.json
\end{verbatim}
を実行した．$3\le Q\le1001$ の奇数五百個について，正 reverse 側の局所模型記録は十五件，
三 core 接着の完成は零件であった．十五件は十四配置への重複割当を含む．
この範囲で外部二標識の交換を除いた正 reverse 側模型は，
$R=(-1,-2,2)$ に対する次の四つである．scale 順は $(S,T,U)$ である．
\begin{center}\small
\begin{tabular}{r|rrr|r}
低 $D$&$S$&$T$&$U$&$Q$\\\hline
2&6&3&6&13\\
9&4&8&2&5\\
9&34&5&29&59\\
9&22&26&38&77
\end{tabular}
\end{center}
この表を全模型の分類として引用してはならない．
$Q>1001$ の未検証範囲を残した有界探索であり，公開稿の無上界排除は正例外型だけである．

\subsubsection{負例外で保持すべき小分母の構造}

負例外の共有高 ray は両側で低 ray と負に結ばれるため，
通常型のように正 reverse 側から二 scale の和を固定する手順をそのまま流用できない．
ただし $Q=3,5$ は強い制約である．公開稿の第二 template
$u=-(Q-1),v=-1$ では，共有高 rows は外部標識 $(0,1)$ の順で
\[
 (X,-(X+Q-1),-1),\qquad (-(Y+1),Y,-(Q-1))
\]
となる．辺 $34$ の puncture $2$ の node は二であるから，
$n_{34;0}+n_{34;1}=3$ であり，
\[
 X-Y=Q(1-n_{34;0})+1,\qquad 0\le n_{34;0}\le3
\]
を得る．第一 template $Q=3,u=-1,v=-2$ では rows は
\[
 (X,-(X+1),-2),\qquad (-(Y+2),Y,-1)
\]
であり，$X-Y=5-3n_{34;0}$ である．
これに低 $D$ と共有高 $D$ の相殺を加え，二つの局所 pair 系を同時に消去することが次の具体的な経路になる．
中心の正の枝は第二 template の $L=3$ だけに残り，
$T=S+2$，$r<S<2r$ を常に保持できる．
負の枝には $S=T=w$ と $w^2<2(r-w)(r-w+1)<4w^2$ を用いる．

\subsubsection{継続上の判断}

今回の進展を前提にするなら，通常型の五つの固定 $(R,D_0)$ 問題，
または負例外の $Q=3$ における共有高 row の連立消去が，次の作業を限定できる具体的な候補である．
$G_{5,14}$ の低辺補 graph $2K_2$ は今回の還元の対象外であり，未着手である．
支持全体に依存する $G_{5,15}$ の定理を，$G_{5,14}$ や $G_{5,16},G_{5,17}$ の部分 graph に適用してはならない．
新しい構造定理や構成例が別方向の価値を高める場合には，この優先順位を変更してよい．

\subsection{Book の低三角形：残る二つの sector と完全復元}
\label{nb20:section}
\label{nb21:section}

$G_{5,13}=\{01,02,03,04,12,13,14\}$ の二価 cores $2,3,4$ を扱う．
正の整数局所 charts，奇数高分母，低辺の node 非負・総和五，
対向低 $D$ 相殺の下で，canonical 非零低低 entry は $-L$，$L\in\{1,3,4\}$ となる．
全正枝，全負枝，二負枝は全て排除されており，残るなら負枝は一つ，正枝は二つである．
全ての低 $D$ は非零で，各 core は一正一負の低 $D$ を持つ．
これらは高中心 $0,1$ の完成を仮定しない必要条件である．

\paragraph{現在残る二つの sector．}
高 ray $0$ へ正 entry を持つ低 row の先を next とすると，有向 node 三角形になる．
canonical $A=0$ を $\sigma=+$，$A=1$ を $\sigma=-$ とする．
前者では next が $C$，後者では next が $D$ である．
next の低 $D$ の共通符号を $\tau$ とし，唯一の負枝 $N$ を $\sigma_N=+$ に正規化すると，
残る可能性は次だけである：
\[
 \mathrm I:\quad N^+\to P_1^+\to P_2^+\to N^+,
 \qquad \tau=+,\qquad (L_{P_1},L_{P_2})=(3,3),
\]
\[
 \mathrm {II}:\quad N^+\to P_1^+\to P_2^-\to N^+,
 \qquad \tau=+,\qquad (L_{P_1},L_{P_2})=(4,3).
\]
負枝は $S_N=T_N=w$，$H=r_N-w$ に対して
\[
 w\ge135,\qquad H\ge w
\]
を満たす．型 II ではさらに
\[
 S_{P_2}=w-1,\qquad T_{P_2}=w+1,
 \qquad n(P_2N)=(0,0,5)
\]
であり，node の順は高 ray $0,1$ と残る低 ray である．
正枝の小さい $L$ に対する符号・分母順序の補題と負枝の符号計算を合わせると，
この正規化の下で三つ全ての二価 core に
\[
 q_{i0}<q_{i1}\qquad(i=2,3,4)
\]
が成り立つ．これは高中心二つに接続する分母三つの厳密な比較である．
中心の四高分母条件への接続は未解決であり，book 自体の不存在とはしていない．

\paragraph{再利用する node offset．}
有向順序で next row を $(a_i,-b_i,-c_i)$，
previous row を $(-d_i,e_i,-f_i)$ と書く．最初の二成分は高 ray $0,1$ に対応する．
\[
 b_i=a_i+1-c_i,\quad d_i=e_i+1-f_i,\quad
 c_if_i=0,\quad c_i+f_i=L_i\in\{1,3,4\},
\]
\[
 (a_i,e_i,c_i,f_i)=
 \begin{cases}
 (r_i,T_i,0,L_i),&\sigma_i=+,\\
 (T_i,r_i,L_i,0),&\sigma_i=-.
 \end{cases}
\]
辺 $i\to i+1$ の node は
\[
 (n_0,n_1,n_{\mathrm{low}})
 =(1-a_i+d_{i+1},\ 1+b_i-e_{i+1},\ 1+c_i+f_{i+1})
\]
であり，非負性は有限 offset
\[
 c_i+f_{i+1}\le4,\qquad e_{i+1}=a_i+t_{i+1},\qquad
 f_{i+1}-2\le t_{i+1}\le2-c_i
\]
と同値である．残余 sector では $\tau=+$ なので，
負の $D$ 端から正の $D$ 端へ向けた周期は node 周期と逆向きになる．
二つの向きは定義が異なるため，計算中に同一視しない．
また高辺 $01$ があるので，負枝を禁じる低 wedge の補題は適用できない．

\paragraph{低係数と局所 sector の比較式．}
一般二高標準形の低係数は，$\delta=T-S$ として
\[
 D_C=\delta-\frac Sr s+\frac T{r+1}z,\qquad D_D=s-z-1.
\]
負枝では $w,H\ge2$，$w^2/2<H(H+1)<2w^2$ から
\[
 \frac{17}{10}w-\frac12<r<\frac{17}{7}w.
\]
$E_-=1-w^2/[H(H+1)]$ とすると，非零低 $D$ の下で
\[
 D_C>0\Longleftrightarrow E_->0\Longrightarrow r\ge2w,
 \qquad
 D_C<0\Longleftrightarrow E_-<0\Longrightarrow r\le2w-1.
\]
局所的に負と零の低 $D$ を持つ唯一の模型は
$(r,s,z,S,T)=(4,4,5,2,2)$，$(D_C,D_D)=(0,-2)$ であり，その低三角形完成は排除済みである．
零低辺を新たに残余 case として加え直す必要はない．

正枝では $\delta\in\{2,3\}$，$h=S-r$，$d=\delta-1$，$k=h+d$ と置く．
\[
 \alpha=\frac{sh}{r},\qquad\beta=\frac{zk}{r+1},\qquad
 F=(2r+1)h-r^2d,
\]
\[
 F\alpha\beta=(\beta-\alpha-1)hk,
 \qquad (rs+1)\beta-((r+1)z+1)\alpha=1.
\]
一正一負の低 $D$ を持つとき，$D_C$ と $F$ は同符号であり，
$D_C<0$ は $q_A>q_B$，$D_C>0$ は $q_A<q_B$ に対応する．
再利用する reverse bounds は
\[
 F>0\Longrightarrow
 T\ge\begin{cases}(3r+4)/2,&\delta=2,\\2r+3,&\delta=3,\end{cases}
 \qquad
 F<0\Longrightarrow
 T\le\begin{cases}\lfloor(3r+3)/2\rfloor,&\delta=2,\\2r+2,&\delta=3.\end{cases}
\]
$F<0$ では $-F\le hk$，従って
\[
 d r^2\le2hr+h(h+\delta).
\]
$\beta=1$ の reverse data は $(r,S,T,z)=(3,4,6,2),(2,3,6,1)$ だけである．
これは Bezout 恒等式から $k(1+2z-dz^2)=\delta$ を得て，
$k\ge\delta$ を使うと分かる．両共鳴族は node offset により負 $w\le5$ を強制するため，
現在の残余 sector では現れない．従って残余正枝の $F<0$ 側では $\beta\ge2$ であり，
\[
 0<\frac{-F}{hk}<\frac12,
 \quad 2d r^2<h^2+4rh+(\delta+1)h,
 \quad h>\frac37r-1,
 \quad T>\frac{10}{7}r+1
\]
を使える．

\paragraph{排除済み sector と有限 certificate の境界．}
二負枝を仮定すると，全ての $\sigma$ と $\tau$ の配置で負 core の一つに
$w\le50$ が導かれる．この解析的範囲の全 scale 復元は完成を排除する．
一負枝の二つの node 型では，上の局所不等式と有限 offset から
\begin{center}
\begin{tabular}{lc}
追加で排除した sector & 導出された負 $w$ の上界\\\hline
型 I，$\tau=-$ & $134$\\
型 I，$\tau=+$，少なくとも一方の正枝が $L=4$ & $8$\\
型 II，$\tau=-$ & $11$\\
型 II，$\tau=+$，$L_{P_1}=3$ & $6$
\end{tabular}
\end{center}
となる．$2\le w\le134$ の全 scale 復元では初期 reverse data が $6392$ 個，
負 canonical cores が $1089$ 個，左右の相手が合計 $169$ 個である．
両隣がともに非空なのは六 core だけで，最後の十一組の $D$ 残差は全て非零だった．
\nolinkurl{verification/verify_book_one_negative_sector_reduction_v21.py} と
対応する JSON が全初期 core・全相手集合・最後の十一組を保存している．
ここでは $w$ の上界自体が仮想的な sector から導かれている．
一方，残余二 sector の $w\ge135$ 側への一様上界はまだない．

\paragraph{固定 core の完全復元への入口．}
一つの二価 core の行列を固定すると，各低辺の二十一 node 分配が
相手の低 row 全体を $p^j_{ik}=1-n_k-p^i_{jk}$ に固定する．
対向 $D$ の相殺は，その低 ray の指定値 $D_*$ も固定する．
canonical reverse data $(r,S,T)$ は有限に復元できる：
固定 row が $C$ なら $r$ が固定され，正枝は $r<S<r\delta$，
負枝は $2\le w<r$ と既知の $H$ 範囲を満たす．
固定 row が $D$ なら $S,T$ が固定され，正枝は $S/\delta<r<S$，
負枝は $w^2/2<H(H+1)<2w^2$ を満たす．

$h=S-r$，$\delta=T-S\in\{0,2,3\}$ とすると，残る scale は
\[
 Fsz+As+Bz+C=0,\qquad \ell s+mz+n=0,
\]
\[
 F=(2r+1)h-r^2(\delta-1),\quad
 A=h(r+1),\quad B=-r(h+\delta-1),\quad C=r(r+1),
\]
\[
 (\ell,m,n)=
 \begin{cases}
 (-S(r+1),rT,(\delta-D_*)r(r+1)),&C\text{ ray},\\
 (1,-1,-D_*-1),&D\text{ ray}
 \end{cases}
\]
を満たす．$\gcd(2r+1,r^2)=1$ と $\delta\in\{0,2,3\}$ から $F\ne0$，
また $\ell m\ne0$ なので，消去した二次式
\[
 -F\ell s^2+(-Fn+Am-B\ell)s+(Cm-Bn)=0
\]
の最高次係数は非零である．固定 reverse data ごとに高々二つの scale を，
整数平方判定と整除から全て復元する．七・十三の共鳴族もこの二次式に含まれる．
二つの低隣接 core の有限集合を作り，残る低辺を検査すれば，全低三角形完成を得る．

負の固定 reverse data $(w,H)$ 自体の全 scale は
\[
 c=\frac{sH}{r},\quad e=\frac{z(H+1)}{r+1},\quad
 E_-=\frac MN\ne0,\quad E_-ce=e-c+1,
\]
\[
 (N-Mc)(N+Me)=N(N-M)>0
\]
から復元する．$c\ge2,e\ge1$ により両因子は正なので，全正約数を取ればよい．
この二つの復元法はいずれも $s,z$ または高分母の cutoff を必要としない．
汎用の固定 core 完成器は \nolinkurl{research/book_low_triangle_completion_v19.py} にあり，
別実装による第21版の整数行列・正の weights・相手復元は
\nolinkurl{verification/verify_book_two_negative_obstruction_v21.py} にある．
低 $D$ を元の pair equation に直接代入する後者は，旧 generator を import しない．

\paragraph{次に比較する量．}
残る二 sector に限り，node offset，低 $D$ の大きさ，
$q_{i0}<q_{i1}$ と中心二 core の四高分母条件を同時に扱う．
型 II の強制 row $S_{P_2}=w-1,T_{P_2}=w+1$ は逆復元の入口になる．
固定 core の有限性を core 自体が動く族の有限性と取り違えず，
新しい一様上界または反例構成を優先する．

% v21: current frontier and reusable next steps; stable support identifiers.
\subsection{残る五型と，再開する具体的な入口}
\label{ns20:support-agenda}

現在の高分母 support の必要候補は\emph{五型・六十六標識}である．
再開時は最初に \nolinkurl{00_CONTINUATION_POLICY_ja.txt} を読む．
以下の順位は，面白さとインパクトを最大化するために柔軟に変更する．
全候補は五頂点を使い，最小次数二以上，二点連結，辺数七以上，$b_1\ge3$ である．
各二価頂点の二つの高近傍はともに次数四である．
二価頂点数は零・一・三に限られ，一個の場合は $G_{5,15}$，
三個の場合および七高辺の場合は book $G_{5,13}$ となる．
表の $ij$ は無向辺，$N$ は標識付き個数であり，補 graph の孤立頂点は省略する．
\begin{center}
\footnotesize
\setlength{\tabcolsep}{4pt}
\begin{tabular}{crclcr}
型&$|E|$&次数列&$E$&低辺の補 graph&$N$\\\hline
$G_{5,13}$&7&$(44222)$&$\{01,02,03,04,12,13,14\}$&$K_3$&10\\
$G_{5,14}$&8&$(43333)$&$K_5-\{14,23\}$&$2K_2$&15\\
$G_{5,15}$&8&$(44332)$&$K_5-\{24,34\}$&$P_3$&30\\
$G_{5,16}$&9&$(44433)$&$K_5-\{34\}$&$K_2$&10\\
$G_{5,17}$&10&$(44444)$&$K_5$&$\varnothing$&1
\end{tabular}
\end{center}
全 $1024$ 標識付き graph の現行検算は
\nolinkurl{verify_support_agenda_v20.py} と
\nolinkurl{support_agenda_verification_v20.json} に保持する．
候補数は整数接着，atlas，幾何的非同型類の存在数ではない．

\paragraph{第一候補：book の残余二領域．}
二価 cores $2,3,4$ の低三角形には，負枝がちょうど一つ必要である．
唯一の負枝を $\sigma_N=+$ と正規化すると，残る型は
\[
\begin{array}{c|c|c}
 & \text{node 周期} & (L_{P_1},L_{P_2})\\ \hline
\mathrm I&N^+\to P_1^+\to P_2^+\to N^+&(3,3)\\
\mathrm {II}&N^+\to P_1^+\to P_2^-\to N^+&(4,3)
\end{array}
\]
であり，双方 $\tau=+$，$w\ge135,H\ge w$ である．
全ての二価 cores で $q_{i0}<q_{i1}$ も必要となる．
型IIでは $S_{P_2}=w-1,T_{P_2}=w+1$ と最後の node $(0,0,5)$ が固定される．
この固定差と分母順序を，全 scale の因子式と有限 node offset に同時に課す．
新しい成長不等式と二つの共鳴族の扱いを保持し，単なる探索上限の拡張は優先しない．
二領域の存在・不存在はともに未確定である．
低三角形が残った場合に四高中心 $0,1$ と $q_{01}$ を復元する．

\paragraph{並行候補：$G_{5,15}$ の五つの固定低データ問題．}
この還元では標識を $H=K_5-\{23,24\}$ と取り直し，二価 core を $2$，
共有高分母を $Q=q_{34}$ とする．表の代表とは標識置換だけが異なる．
正例外型は全排除され，通常型 $v=-u$ は四中心・十四配置へ限られる．
負例外型は $Q=3,5$ の二 template となる．
通常型は各固定 $Q$ について cores $2,3,4$ と相互三辺を完全に有限復元できるが，
$Q$ の一様上界を与えるものではない．
次は正 reverse 側の固定低 row と要求 $D_0$ の五問題
\[
((-1,-4,4),5),\quad ((-1,-1,1),-9),\quad
((-1,-2,2),2),\quad ((-1,-2,2),-7),\quad ((-1,-2,2),9)
\]
を無上界で分類する．外部二標識の交換は許している．
負例外では $Q$ が小さくても二高中心の自由度が残るため，別途全尺度を扱う．
$G_{5,14}$ と $G_{5,16},G_{5,17}$ も必要候補として保持する．
四高側へ三高 matching 定理を無条件に拡張しない．

\paragraph{三高 $D$ 問題：単一合同式と残余式の交わり．}
\ref{nhd21:reconstruction} の五データ復元を用いる．
整数尺度には $m=F^2+F+1\mid FE-J$ が必要十分であり，
$N=E^2+EJ+J^2>0$ と $m\mid N$ により各固定 $(s,t,x,y,z)$ の $b$ は有限となる．
ただしノルム条件だけでは合同式に足りず，尺度を復元しても元の pair 残余式を残す．
$j=t-y$ を正と仮定してはいけない．
五変数自体への一様上界，または残余式との無上界矛盾を目標とする．
正整数候補が得られたときは明示的な正 weights まで復元する．

\paragraph{再利用と，解決済みの型．}
$K_4$，二葉三角形，$G_{5,5},G_{5,7},G_{5,9},G_{5,10},G_{5,11},G_{5,12}$ は再探索しない．
reverse data の解析的上界，点・全半直線表示，分母一致と $D$ の奇偶障害は再利用できる．
ただし相手の高次数と例外接着の仮定を毎回確認する．
$G_{5,15},G_{5,16},G_{5,17}$ が $K_4$ を部分 graph に含むだけでは，
support 全体が $K_4$ の場合の排除を移せない．
高分母 support $H$ と素数 support $G_\ell$ を区別する．
dense mod 五の二模型では少なくとも二辺の五進付値一を保つ．
整数候補が残れば actual weights，basin areas，integer framings，作用切片へ進む．
幾何的同値，別退化，局所置換への転換も，研究全体の価値に応じて判断する．



\subsection{第三 corner の有限条件と一辺で整合する無限族}

二つの cores $i,j$ を共有辺 $ij$ で比較し，puncture $k$ に対して
\[
 A=q_{ik},\quad B=q_{jk},\quad r=p^i_{kj},\quad s=p^j_{ki}
\]
が既知とする．第三 core $k$ の二つの numerators は
$e=Au-r$，$f=Bv-s$ と書かれ，$u=1-n_{ik;j}$，$v=1-n_{jk;i}$ である．
各共有直線の全 node 数五と非負性を使うと，必要条件は
\[
 (Au-r)(Bv-s)=AB-1,
 \qquad u,v\in\{-4,-3,-2,-1,0,1\}
\]
となる．この三十六候補の検査は分母の大きさによらず使える．
三 punctures 全てで通っても，第三 core の全データや actual weights の同時実現はまだ必要である．

再利用できる局所無限族として，$n\geq1$，$Q_n=n^2+n+1$，$R_n=n^2-n+1$ に対し
\[
 q=(Q_n,1,1,R_n),\qquad
 P_n=\begin{pmatrix}
 0&-(n+1)&1&-(n^2+1)\\
 -n&0&0&n-1\\
 n(n+1)&-(2n+1)&0&-n(n-1)\\
 -n^2&n&-1&0
 \end{pmatrix},
\]
\[
 w=(Q_n-1/2,1,1,n-1+1/(2n)),\qquad
 D=(-Q_n,2n+1,-1,R_n)
\]
を持つ正局所模型 $L_n$ がある．$L_n$ の $Q_n$ 端と $L_{n+1}$ の $R_{n+1}=Q_n$ 端では
$D$ が相殺し，outgoing および incoming numerators の和は，適切な共通順序で
$(0,0,-2Q_n)$ となる．node 配分は $(1,1,3)$，star shears は $(0,0,2)$ である．
しかし $n\geq2$ では residue matching が一意で，第三 corner は
\[
 (R_nu+n^2)\bigl(Q_{n+1}v+n^2+2n+2\bigr)
   =n^2(n+1)^2+2
\]
を要求し，六値の $u,v$ の全場合が一般の多項式不等式で排除される．
この族は局所模型と単独辺の整数整合の存在例であり，actual weights や切片まで接続した大域例ではない．

\subsection{既存の有界排除の量化と，残る研究の入口}

$q_{ij}\le101$ の全範囲では全 $q_{ij}=1$ が証明済みである．
奇数性から，高分母が存在すれば最大値は百三以上となる．
この結論は分母が全て百一以下という仮定を必要とし，現在の五候補の無上界排除を意味しない．

完全性に必要な手法を残す．二高分母以上の core では $q_0,q_3>1$ と並べ替え，
$p_{01},p_{02},p_{13}$ をそれぞれ $q_0q_1-1,q_0q_2-1,q_1q_3-1$ の signed divisors から選ぶ．
残りは row sums と pair products で決まるため，numerators に独立の探索上限は不要である．
一高分母では完全な高 ray の標準形を保持する．端点の fingerprint
\[
 \left(q,D,\operatorname{sort}(p_1,p_2,p_3\bmod q)\right)
\]
について，対向端での $D$ と residues の符号反転を要求し，相手のない候補を反復削除する．
これは必要条件による削除なので，大域接着の十分条件とはしない．
元の完全 catalog と証明書は
\nolinkurl{verify_quintic_denominators_v10.py} と
\nolinkurl{quintic_denominators_verification_v10.json} に保持し，版番号を付け替えない．

次の対象は $G_{5,13},G_{5,14},G_{5,15},G_{5,16},G_{5,17}$ である．
book では全正枝・全負枝・二負枝を除き，一負枝の残る二つの sector だけを調べる．
全低 $D$ は非零であり，各 core で一正一負となって有向三角形を作る．
有限な node offset と固定 core 完成法を，この符号周期の条件と同時に使う．
三高分母の同時非負性では非特異な canonical Schur 系に集中できるが，
一般の局所模型の分類には正整数解を持つ特異系を保持する．
必要な整数候補が得られれば全 core の weights，整数 framings，実作用切片まで検討する．
密な mod 五模型の整数 lifting では，五進付値一の辺が少なくとも二本ある条件を保持する．

$G_{5,5},G_{5,7},G_{5,9},G_{5,10},G_{5,11},G_{5,12}$ は全排除済みであり，再探索しない．
その証明から，共通低 ray の $C/D$ 型，matching が強制する符号，
未知分母を切らない点・半直線の処理を取り出して次の候補へ使う．
等二高分母の負和，素数 support の合同障害，または部分的に整合する整数模型を，
任意分母の全体に対する不存在・存在へ無条件に移さない．
